\pdfoutput=1
\documentclass[11pt]{amsart}
\usepackage{amsmath,amssymb,amsthm,mathtools}
\usepackage[margin=1.1in]{geometry}
\usepackage{booktabs,array,enumitem}
\usepackage{xcolor}
\usepackage[expansion=false]{microtype}
\usepackage[colorlinks=true,linkcolor=blue!60!black,citecolor=green!40!black,urlcolor=blue!60!black]{hyperref}

\numberwithin{equation}{section}
\newtheorem{theorem}{Theorem}[section]
\newtheorem{proposition}[theorem]{Proposition}
\newtheorem{lemma}[theorem]{Lemma}
\newtheorem{corollary}[theorem]{Corollary}
\theoremstyle{definition}
\newtheorem{definition}[theorem]{Definition}
\theoremstyle{remark}
\newtheorem{remark}[theorem]{Remark}

% conjectures of Berndt et al., with their own numbers
\newtheorem*{bconj@inner}{Conjecture \bconjlabel}
\makeatletter
\newenvironment{bconj}[2][]{\def\bconjlabel{#2}%
  \if\relax\detokenize{#1}\relax\begin{bconj@inner}\else\begin{bconj@inner}[#1]\fi}{\end{bconj@inner}}
\makeatother
\newcommand{\bconjlabel}{}

\newcommand{\Z}{\mathbb{Z}}
\newcommand{\N}{\mathbb{N}}
\newcommand{\R}{\mathbb{R}}
\newcommand{\C}{\mathbb{C}}
\newcommand{\E}{\mathbb{E}}
\DeclareMathOperator{\sgn}{sgn}
\DeclareMathOperator{\sq}{sq}
\newcommand{\dpar}[1]{(\!(#1)\!)}
\newcommand{\Tpart}{\mathcal{T}}
\newcommand{\kz}{\kappa_0}
\makeatletter\g@addto@macro\UrlBreaks{\do\_\do\.}\makeatother
\DeclareUrlCommand\lean{\urlstyle{tt}}
\setlength{\emergencystretch}{1.5em}

\title[A theta-function sum and a conjecture of Berndt et al.]{An asymptotic formula for a theta-function sum and a conjecture of Berndt, Bhat, Meyer, Xie and Zaharescu}
\author{Samson Fraser}
\email{samsonmfraser@gmail.com}
\address{Independent researcher}
\subjclass[2020]{Primary 11L03; Secondary 11L05, 11F20, 11N25}
\keywords{theta function, Hardy sums, cotangent sums, short Kloosterman sums, divisors in arithmetic progressions}
\date{}

\begin{document}
\begin{abstract}
For an odd prime $p$ let $S(p)=\sum_{h,j=1}^{p-1}(-1)^{j+1+\lfloor hj/p\rfloor}$, a sum arising in the transformation formula of the theta function $\vartheta_3$. Berndt, Bhat, Meyer, Xie and Zaharescu conjectured that for any fixed $C_1,C_2,\varepsilon>0$ one has $C_1p<S(p)<C_2p^{1+\varepsilon}$ for all sufficiently large primes $p$. We prove this conjecture by showing that
\[
S(p)=\tfrac12\,p\log p+O\bigl(p(\log p)^{19/20}\bigr).
\]
The proof starts from an exact formula, obtained from a reciprocity law for half-range cotangent sums, that expresses $S(p)$ through tangent sums $\sum_t\tan(\pi pt/k)$. Large moduli $k$ are treated with Kloosterman's bound. For smaller moduli each tangent sum is split into a head and a tail: the heads are bounded on average over $k$ by the Bourgain--Garaev bound for short Kloosterman sums, and the tails by a new elementary estimate for the large divisors of the terms of an arithmetic progression. A formalization in the Lean proof assistant derives the conjecture from the Bourgain--Garaev bound in the range used here; that bound enters the formalization as a hypothesis.
\end{abstract}
\maketitle

\section{Introduction}\label{sec:intro}

\subsection{The sum and the conjectures}
For integers $h$ and $k\ge1$ let
\begin{equation}\label{eq:Shk}
S(h,k)=\sum_{j=1}^{k-1}(-1)^{j+1+\lfloor hj/k\rfloor},\qquad S(k)=\sum_{h=1}^{k-1}S(h,k).
\end{equation}
The sum $S(h,k)$ is one of the arithmetic sums attached to theta functions that were introduced by Hardy~\cite{b14} and studied by Berndt~\cite{b2} and by Berndt and Goldberg~\cite{b3}. It occurs in the transformation formula of $\vartheta_3(z)=\sum_{n\in\Z}e^{\pi in^2z}$: if $\bigl(\begin{smallmatrix}a&b\\c&d\end{smallmatrix}\bigr)$ lies in the theta group and $c>0$, then under $z\mapsto(az+b)/(cz+d)$ the function $\log\vartheta_3$ changes by elementary terms and the term $\frac{\pi i}4S(d,c)$ \cite[Theorem~4.1]{b2}, much as the Dedekind sum occurs in the transformation formula of the Dedekind eta function.\footnote{In~\cite{b2} the sum runs over $1\le j\le c$, which adds $1$ to $S(d,c)$ when $c+d$ is odd; see also \cite[Section~2]{b1}.}

Berndt, Bhat, Meyer, Xie and Zaharescu~\cite{b1} studied the complete sums $S(k)$ for prime $k$. For an odd prime $p$,
\[
S(p)=\sum_{h=1}^{p-1}\sum_{j=1}^{p-1}(-1)^{j+1+\lfloor hj/p\rfloor}
\]
is a sum of $(p-1)^2$ signs. A few values:
\begin{center}\small
\begin{tabular}{l*{10}{r}}
\toprule
$p$ & 3 & 5 & 7 & 11 & 13 & 101 & 233 & 1009 & 3119 & 999983\\
$S(p)$ & 2 & 4 & 10 & 14 & 20 & 268 & 440 & 4408 & 9194 & 6306294\\
\bottomrule
\end{tabular}
\end{center}
Independent random signs would give a sum of size about $p$ and of either sign. Instead $S(p)$ is positive in all computed cases and grows slightly faster than $p$. On the basis of such data the authors of~\cite{b1} made the following conjectures, in which $k$ denotes a prime.\footnote{We use the numbering of the published version of~\cite{b1}, which agrees with arXiv versions 2 and~3; the formal statements in~\cite{b11} follow the first arXiv version and number them differently (Appendix~\ref{app:lean}).}

\begin{bconj}[positivity]{1.1}
$S(k)>0$ for every prime $k>5$.
\end{bconj}
\begin{bconj}{3.1}
$S(k)>2k$ for every prime $k>233$.
\end{bconj}
\begin{bconj}{3.2}
$S(k)>3k$ for every prime $k>3119$.
\end{bconj}
\begin{bconj}[growth]{3.3}
For any fixed $C_1,C_2,\varepsilon>0$ one has $C_1k<S(k)<C_2k^{1+\varepsilon}$ for all sufficiently large primes $k$.
\end{bconj}

They write that ``Conjecture~1.1 appears to be very difficult to prove'', and trace the difficulty to ``our lack of knowledge about the distribution of the residues of $hj$ modulo $k$'', that is, to the distribution of the points on the modular hyperbolas $hj\equiv c\pmod k$ (see Shparlinski's survey~\cite{b25}).

\subsection{Results}
We determine the size of $S(p)$ and prove Conjecture~3.3.

\begin{theorem}\label{thm:main}
\textup{(a)} For every odd prime $p$,
\[
|S(p)|\le\frac2{\pi^2}\,p\,(\log p+2)^2.
\]
\textup{(b)} As $p\to\infty$ over the primes,
\[
S(p)=\frac p2\log p+O\bigl(p(\log p)^{19/20}\bigr).
\]
In particular, Conjecture~\textup{3.3} holds, and for every fixed $B$ we have $S(p)>Bp$ for all sufficiently large primes $p$.
\end{theorem}

So the $(p-1)^2$ signs cancel almost completely, and what is left is a bias of order $\log p$ times the size of a random sum. Part~(a) is elementary. Part~(b) is the substance of the paper; it implies $S(p)>0$, $S(p)>2p$ and $S(p)>3p$ for all sufficiently large $p$. The threshold depends on the implied constant in the theorem of Bourgain and Garaev used in Section~\ref{sec:kloosterman}, which has not been computed, and on the largeness conditions in the proof, which as written require at least $\log\log p>4073$ (Remark~\ref{rem:optimality}(c)). So part~(b) says nothing about the primes within computational reach. For small primes we have the following computation (Section~\ref{sec:computations}).

\begin{proposition}\label{prop:numerics}
For all odd primes $p\le6\cdot10^6$ we have $S(p)>0$; $S(p)>p$ if $p>5$; $S(p)>2p$ if $p>233$; $S(p)>3p$ if $p>3119$; $S(p)>4p$ if $p>58013$; and $S(p)>5p$ if $p>869399$. None of these thresholds can be lowered.
\end{proposition}

So Conjectures~1.1, 3.1 and~3.2 hold for $p\le6\cdot10^6$ and for all sufficiently large $p$; the range in between remains open.

The new ingredient of the proof is an elementary estimate for the large divisors of the terms of an arithmetic progression. Divisors are positive: for a nonzero integer $n$, $d\mid n$ means that $d$ is a positive divisor of $|n|$.

\begin{theorem}\label{thm:divisors}
Let $0<\gamma_0<1$. There is $L_0(\gamma_0)$ such that the following holds for every integer $q$ with $\log q\ge L_0(\gamma_0)$. Put $T=\exp((\log q)^{\gamma_0})$. Let $a$ be an integer coprime to $q$, let $R\ge1$, and let $\mathcal I$ be a set of at most $R$ consecutive integers such that $0<|a+qt|\le q^2$ for all $t\in\mathcal I$. Then
\[
\sum_{t\in\mathcal I}\ \sum_{\substack{d\mid a+qt\\ d>TR}}\frac1d\;\le\;T^{-\gamma_0/4}.
\]
\end{theorem}

For $R$ random integers of size at most $q^2$ the left-hand side would be about $1/T$. The point of Theorem~\ref{thm:divisors} is that the length $R$ of the progression is arbitrary: bounds for divisor sums along progressions, such as Shiu's theorem~\cite{b24}, require the length to be at least a fixed power of the modulus, whereas we need every length from $1$ to about $q^{3/4}$.

\subsection{The method}
\emph{An exact formula.} A finite Fourier expansion modulo $2p$ turns the sum of signs into a sum of products of cotangents (Lemma~\ref{lem:dual}):
\[
S(p)=\frac2p\sum_{\substack{0<k<p\\ k\text{ odd}}}\cot\frac{\pi k}{2p}\,C_p(k),\qquad C_p(k)=\sum_{\substack{0<a<p\\ a\text{ odd}}}\cot\frac{\pi ka}{2p}.
\]
Since $|C_p(k)|\le C_p(1)\le\frac p\pi(\log p+2)$, this gives Theorem~\ref{thm:main}(a). The half-range cotangent sums satisfy a reciprocity law (Lemma~\ref{lem:reciprocity}),
\[
C_p(k)=\frac pk\,F_k(p)+I\Bigl(\frac kp\Bigr),\qquad F_k(n)=\sum_{t=1}^{(k-1)/2}\tan\frac{\pi nt}k,
\]
with an explicit integral $I$ satisfying $\pi\lambda I(\lambda)=\log\frac1\lambda+\kz+O(\lambda^2)$. The terms $I(k/p)$ produce the main term $\frac p2\log p+c_0p$, with an explicit constant $c_0=0.2399\ldots$. For $k\ge K_W=p^{3/4}(\log p)^4$ we use cancellation between different $k$: summed over an interval of $k$, the $C_p(k)$ become incomplete Kloosterman sums, and Kloosterman's bound shows that these moduli contribute $O(p/\log p)$. The result (Theorem~\ref{thm:decomp}) is
\begin{equation}\label{eq:intro-decomp}
S(p)=\frac p2\log p+c_0p+\frac{4p}\pi\sum_{\substack{3\le k<K_W\\ k\text{ odd}}}\frac{F_k(p)}{k^2}+O\Bigl(\frac p{\log p}\Bigr).
\end{equation}
Since $|F_k(p)|\le\frac k\pi(\log k+2)$, the sum over $k$ is trivially $O((\log p)^2)$. The problem is to show that it is $o(\log p)$, for every large prime $p$.

\emph{The tangent sums.} The quotient $F_k(p)/k$ depends only on $\theta_k=u_k/k$, where $2pu_k\equiv1\pmod k$, and it is a weighted sum of signs along the multiples of $\theta_k$ (Lemma~\ref{lem:sign-form-tangent}):
\[
\frac{F_k(p)}k=-\sum_{\substack{0<n<k\\ n\text{ odd}}}\frac1k\cot\frac{\pi n}{2k}\,\sq(n\theta_k),\qquad \sq(y)=\sgn\bigl(\tfrac12-\{y\}\bigr).
\]
It is large when $\theta_k$ is close to a fraction with small denominator; for instance $F_k(p)/k=-\frac1\pi\log k+O(1)$ whenever $k\mid p+1$. Such moduli occur for infinitely many $p$, so the estimate has to be uniform in $p$. We split the sum over $n$ into a \emph{head} $n\le N$ and a \emph{tail} $n>N$, where $\log N\asymp(\log p)^{1/20}$.

\emph{Heads.} Each head is $O(\log N)$. For $k<K_*=\exp((\log p)^{9/10})$ we use this bound; with the weights $1/k$ these moduli contribute $O(\log N\log K_*)=O((\log p)^{19/20})$, which is the error term in Theorem~\ref{thm:main}(b). For $K_*\le k<K_W$ we average the heads over $k$ in dyadic blocks. Expanding the sign function into a trigonometric polynomial turns them into short Kloosterman sums $\sum_{n\le N'}e(bn^*/p)$, where $n^*$ is the inverse of $n$ modulo $p$ and $\log N'\ge\frac12(\log p)^{9/10}$. Here we use the bound of Bourgain and Garaev~\cite[Theorem~16]{b8}, which saves a power of $\log p$ for such sums (Section~\ref{sec:kloosterman}).

\emph{Tails.} The tail is small when the points $n\theta_k$ are well distributed modulo~$1$. An Erd\H{o}s--Tur\'an type argument shows that this holds unless $\|2m\theta_k\|<1/T$ for some $m\le\lceil(\log p)^2\rceil$, where $T=\exp((\log p)^{1/20})$. No estimate for exponential sums is available at this scale, so we count the exceptional moduli instead. Such a $k$ has a \emph{witness}: an integer $r$ with $0<|r|<k/T$ and $k\mid pr-m$. So $k$ is a divisor larger than $T|r|$ of a term of one of the progressions $\{pt\mp m\}$, and Theorem~\ref{thm:divisors} shows that the total $1/k$-mass of these moduli is negligible.

\subsection{The external input and the formal verification}
Apart from standard tools (Vaaler's approximation of the sawtooth function and the Rosser--Schoenfeld bound for $\prod_{\ell\le y}(1-1/\ell)^{-1}$), the proof uses one external result: Theorem~16 of Bourgain and Garaev~\cite{b8}, for sums of length $N$ with $\log N\ge\frac12(\log p)^{9/10}$. This lies far inside the range of $N$ covered by its proof, which is $\log N\ge(\log p)^{2/3}(\log\log p)^{O(1)}$ (Lemma~\ref{lem:bg}; Garaev and Shparlinski~\cite[(1.2)]{b12} state it as $\log N\ge(\log p)^{2/3}(\log\log p)^{c}$ for any fixed $c>3$). For $N\ge p^{4/7}$, which occurs here, the proof of~\cite[Theorem~16]{b8} rests on Weil's bound for Kloosterman sums, so Weil's bound enters the paper through this input. Everything else is proved here, including Kloosterman's bound for complete sums (Lemma~\ref{lem:kloosterman}), which is all that the decomposition of Section~\ref{sec:decomposition} needs.

The deduction of Conjecture~3.3 from this input has been verified in the Lean proof assistant. Starting from the definition of $S$ in the formal-conjectures repository~\cite{b11}, the formalization proves the implication ``Lemma~\ref{lem:bg} $\Rightarrow$ Conjecture~3.3''; the bound of Lemma~\ref{lem:bg} enters as a hypothesis and is not formalized, and every other step is machine-checked. The formal proof follows the route of Sections~\ref{sec:notation}--\ref{sec:proof}, with minor differences; it establishes $S(p)/p\to\infty$ rather than the explicit exponent $\frac{19}{20}$ (Appendix~\ref{app:lean}). The code, the data and the Lean files are available at \url{https://github.com/samsonmaximus/theta-sum}.

\subsection{Related work}
Hardy sums and their reciprocity laws were developed by Berndt~\cite{b2}, Berndt and Goldberg~\cite{b3} and Sitaramachandrarao~\cite{b26}, and L\"ageler~\cite{b18} related $S(d,c)$ to even continued fractions. Reciprocity formulas, bounds and limiting distributions for cotangent sums were obtained by Bettin and Conrey~\cite{b5}, Bettin~\cite{b4} and Bettin and Drappeau~\cite{b6,b7}. These results concern averages over complete residue systems or over all denominators, whereas $S(p)$ involves essentially one point for each modulus $k$. For short Kloosterman sums see Bourgain and Garaev~\cite{b8,b9}, Korolev~\cite{b17} and Garaev and Shparlinski~\cite{b12}.

\subsection{Organization}
Sections~\ref{sec:notation}--\ref{sec:decomposition} prove the cotangent formula and Theorem~\ref{thm:main}(a), the reciprocity law, and the decomposition~\eqref{eq:intro-decomp}. Sections~\ref{sec:kloosterman} and~\ref{sec:tangent} state the input on short Kloosterman sums and treat heads and tails, and Section~\ref{sec:divisors} proves Theorem~\ref{thm:divisors} and can be read on its own. Section~\ref{sec:proof} proves Theorem~\ref{thm:main}(b), Section~\ref{sec:computations} describes the computations, and Appendix~\ref{app:lean} the formal verification.

\section{Notation and the cotangent formula}\label{sec:notation}

This section fixes the notation and expands the sign $(-1)^{\lfloor n/p\rfloor}$ in a finite Fourier series modulo $2p$, which gives the cotangent formula~\eqref{eq:dual} for $S(p)$ (Lemma~\ref{lem:dual}), the starting point of Section~\ref{sec:decomposition}, and the sign form of the sums $C_p(k)$ (Lemma~\ref{lem:sign-form}), used in Proposition~\ref{prop:upper} and Lemma~\ref{lem:Cp-interval}. The same two lemmas give the upper bound $|S(p)|\le\frac2{\pi^2}p(\log p+2)^2$ of Proposition~\ref{prop:upper}, which is part~(a) of Theorem~\ref{thm:main} and is used in Section~\ref{sec:proof}. Nothing in this section depends on the later ones apart from the definition of $S(p)$ in Section~\ref{sec:intro}.

\subsection{Notation}
Throughout, $p$ is an odd prime. We write $e(x)=e^{2\pi ix}$, $\|x\|$ for the distance from $x$ to the nearest integer, $\{x\}=x-\lfloor x\rfloor$, and $\dpar{x}$ for the sawtooth function: $\dpar{x}=\{x\}-\frac12$ for $x\notin\Z$ and $\dpar{x}=0$ for $x\in\Z$. The symbols $\ll$ and $O(\cdot)$ have their usual meaning; implied constants are absolute unless a dependence is indicated by a subscript.

\emph{Signs.} For $m\in\Z$ with $p\nmid m$ we put
\[
\sigma(m)=(-1)^{\lfloor m/p\rfloor};
\]
thus $\sigma(m)=1$ if $m\bmod2p\in(0,p)$ and $\sigma(m)=-1$ if $m\bmod2p\in(p,2p)$, and $\sigma$ is $2p$-periodic and odd. For $y\in\R\setminus\frac12\Z$ we put
\[
\sq(y)=\sgn\Bigl(\frac12-\{y\}\Bigr).
\]
Then $\sigma(m)=\sq(m/2p)$ for $p\nmid m$, since $m/2p\in\frac12\Z$ if and only if $p\mid m$.

\emph{Inverses.} For $k$ coprime to $2p$, $\overline k$ denotes an integer with $k\overline k\equiv1\pmod{2p}$; it is odd. For $n$ coprime to $p$, $n^*$ denotes an inverse of $n$ modulo $p$. If $k$ is odd and $p\nmid k$, then $k\overline k\equiv1\pmod{2p}$ implies $k\overline k\equiv1\pmod p$, hence
\[
\overline k\equiv k^*\pmod p.
\]

\emph{Sums.} For coprime odd positive integers $m,k$ let
\[
C_m(k)=\sum_{\substack{0<a<m\\ a\text{ odd}}}\cot\frac{\pi ka}{2m}.
\]
For odd $k\ge3$ and $n\in\Z$ the tangent sums are
\begin{equation}\label{eq:Fk}
F_k(n)=\sum_{t=1}^{(k-1)/2}\tan\frac{\pi nt}k,
\end{equation}
and $F_1=0$. No term of~\eqref{eq:Fk} is singular: $nt/k\equiv\frac12\pmod1$ would mean $2nt\equiv k\pmod{2k}$, which is impossible because $k$ is odd. For $K>3$ we put
\[
R_K(n)=\sum_{\substack{3\le k<K\\ k\text{ odd}}}\frac{F_k(n)}{k^2}.
\]
The function
\begin{equation}\label{eq:I}
I(\lambda)=\frac1\pi\int_0^\infty\frac{\cosh v-\cosh\lambda v}{\sinh v\,\sinh\lambda v}\,dv\qquad(\lambda>0)
\end{equation}
and the constants
\begin{equation}\label{eq:K0}
\kz=\gamma+\log\frac4\pi=0.8187801401\ldots,
\end{equation}
where $\gamma$ is Euler's constant, and
\begin{equation}\label{eq:c0}
c_0=\frac{\kz}2+\frac{\log2}6+\frac3{\pi^2}\zeta'(2)=0.2399341036\ldots
\end{equation}
enter the reciprocity law and the main term of $S(p)$ (Sections~\ref{sec:reciprocity} and~\ref{sec:decomposition}).

\emph{The function $h$.} Let $h(x)=\frac{\pi x}2\cot\frac{\pi x}2$ for $0<x\le1$ and $h(0)=1$. Then $h$ decreases from $1$ to $0$ on $[0,1]$, is continuously differentiable there, and
\[
0\le1-h(x)\le x^2\qquad(0\le x\le1).
\]
Indeed, $1-y\cot y=\sum_{n\ge1}2\zeta(2n)(y/\pi)^{2n}$ for $|y|<\pi$ has positive Taylor coefficients, so $(1-h(x))/x^2$ is nonnegative and increasing on $(0,1]$, and its value at $x=1$ is $1-h(1)=1$.

\subsection{The cotangent formula and the upper bound}
Recall that $S(p)=\sum_{h=1}^{p-1}S(h,p)$, with $S(h,k)$ as in~\eqref{eq:Shk}. The first lemma uses the reflection $b\mapsto p-b$ to show that the terms with even $a$ cancel.

\begin{lemma}\label{lem:elementary}
For every odd prime $p$ and $1\le a\le p-1$ put $U(a)=\sum_{b=1}^{p-1}(-1)^{\lfloor ab/p\rfloor}$. Then
\begin{enumerate}[label=\upshape(\alph*),leftmargin=*,nosep]
\item $S(p)=\sum_{a,b=1}^{p-1}(-1)^{\lfloor ab/p\rfloor}=\sum_{a,b=1}^{p-1}\sigma(ab)$;
\item $U(a)=0$ for even $a$, so that $S(p)=\sum_{0<a<p,\ a\text{ odd}}U(a)$.
\end{enumerate}
\end{lemma}

\begin{proof}
For $1\le a,j\le p-1$ we have $p\nmid aj$, hence $\lfloor(p-a)j/p\rfloor=j-1-\lfloor aj/p\rfloor$ and
\[
j+1+\Bigl\lfloor\frac{(p-a)j}p\Bigr\rfloor\equiv\Bigl\lfloor\frac{aj}p\Bigr\rfloor\pmod2.
\]
Thus $S(p-a,p)=U(a)$, and summing over $a$ gives the first equality in~(a); the second is the definition of $\sigma$. If $a$ is even, then $\lfloor a(p-b)/p\rfloor=a-1-\lfloor ab/p\rfloor\equiv1+\lfloor ab/p\rfloor\pmod2$, so the terms $b$ and $p-b$ of $U(a)$ cancel; since $p$ is odd, $b\mapsto p-b$ has no fixed point. This proves $U(a)=0$, and the last claim of~(b) follows from~(a).
\end{proof}

\begin{lemma}[Dual formula]\label{lem:dual}
For odd $a$ with $0<a<p$,
\[
U(a)=\frac2p\sum_{\substack{0<k<p\\ k\text{ odd}}}\cot\frac{\pi k}{2p}\cot\frac{\pi ka}{2p},
\]
and consequently
\begin{equation}\label{eq:dual}
S(p)=\frac2p\sum_{\substack{0<k<p\\ k\text{ odd}}}\cot\frac{\pi k}{2p}\,C_p(k).
\end{equation}
\end{lemma}

\begin{proof}
The idea is to expand $f(n)=(-1)^{\lfloor n/p\rfloor}$ in a Fourier series modulo $2p$, whose coefficients are cotangents, and to sum the resulting geometric series in $b$.

\emph{Step 1: the Fourier coefficients.} The function $f$ has period $2p$ and satisfies $f(n+p)=-f(n)$. Hence its Fourier coefficients $\hat f(\nu)=\frac1{2p}\sum_{n=0}^{2p-1}f(n)e(-\nu n/2p)$ vanish for even $\nu$, and for odd $\nu$, with $\zeta=e(-\nu/2p)$ (so that $\zeta^p=-1$),
\begin{equation}\label{eq:fourier-coeff}
\hat f(\nu)=\frac1{2p}(1-\zeta^p)\sum_{n=0}^{p-1}\zeta^n=\frac2{p(1-\zeta)}=\frac1p\Bigl(1-i\cot\frac{\pi\nu}{2p}\Bigr);
\end{equation}
for $\nu=p$ this equals $1/p$.

\emph{Step 2: the sum over $b$.} Let $a$ be odd, $0<a<p$. Then $U(a)=\sum_\nu\hat f(\nu)\sum_{b=1}^{p-1}e(\nu ab/2p)$, where $\nu$ runs over the odd residues modulo $2p$. For $\nu=p$ the inner sum is $\sum_{b=1}^{p-1}(-1)^{ab}=\sum_{b=1}^{p-1}(-1)^b=0$. For odd $\nu\ne p$ put $\vartheta=\nu a/2p$; then $e(p\vartheta)=-1$ because $\nu a$ is odd, $e(\vartheta)\ne1$ because $p\nmid\nu a$, and
\[
\sum_{b=1}^{p-1}e(b\vartheta)=\frac{1-e(p\vartheta)}{1-e(\vartheta)}-1=\frac{1+e(\vartheta)}{1-e(\vartheta)}=i\cot(\pi\vartheta).
\]
Hence
\[
U(a)=\frac1p\sum_{\substack{0<\nu<2p,\ \nu\ne p\\ \nu\text{ odd}}}\Bigl(\cot\frac{\pi\nu}{2p}\cot\frac{\pi\nu a}{2p}+i\cot\frac{\pi\nu a}{2p}\Bigr).
\]

\emph{Step 3: symmetry.} Under $\nu\mapsto2p-\nu$ both $\cot(\pi\nu/2p)$ and $\cot(\pi\nu a/2p)$ change sign (the second because $a$ is an integer). So the imaginary part of the sum vanishes, while the product of the two cotangents is invariant. The terms $\nu$ and $2p-\nu$ are thus equal, and exactly one of them lies in $(0,p)$; this gives the formula for $U(a)$. Equation~\eqref{eq:dual} follows by summing over odd $a$ and using Lemma~\ref{lem:elementary}(b).
\end{proof}

\begin{lemma}[Sign form]\label{lem:sign-form}
Let $k$ be odd with $p\nmid k$. Then
\begin{equation}\label{eq:sign-form}
C_p(k)=\sum_{\substack{0<n<p\\ n\text{ odd}}}\sigma(n\overline k)\cot\frac{\pi n}{2p}.
\end{equation}
In particular $|C_p(k)|\le C_p(1)$.
\end{lemma}

\begin{proof}
For odd $n$ with $0<n<p$ let $w_n\in(0,2p)$ be the residue of $n\overline k$ modulo $2p$; it is odd, and it is different from $p$ because $\overline k$ is prime to $2p$ and $p\nmid n$. Put $a_n=w_n$ if $w_n<p$ and $a_n=2p-w_n$ otherwise. Then $a_n$ is odd, $0<a_n<p$, and $ka_n\equiv\sigma(n\overline k)\,n\pmod{2p}$, so that $\cot(\pi ka_n/2p)=\sigma(n\overline k)\cot(\pi n/2p)$. The map $n\mapsto a_n$ is a bijection of the set of odd integers in $(0,p)$ onto itself; its inverse sends $a$ to the element of $\{ka\bmod2p,\ 2p-(ka\bmod2p)\}$ that lies in $(0,p)$. Summing over $n$ gives~\eqref{eq:sign-form}, and the bound follows since $\cot(\pi n/2p)>0$ for $0<n<p$.
\end{proof}

Combining~\eqref{eq:dual} and~\eqref{eq:sign-form} gives the symmetric formula
\[
S(p)=\frac2p\sum_{\substack{0<k,n<p\\ k,n\text{ odd}}}\cot\frac{\pi k}{2p}\,\cot\frac{\pi n}{2p}\,\sigma(n\overline k).
\]

\begin{lemma}\label{lem:Cm1-elementary}
For odd $m\ge3$, $C_m(1)\le\frac m\pi(\log m+2)$.
\end{lemma}

\begin{proof}
We have $\cot y\le1/y$ for $0<y\le\pi/2$, so $\cot\frac{\pi a}{2m}\le\frac{2m}{\pi a}$ for odd $a<m$. Also $\sum_{a\text{ odd},\,3\le a\le m}1/a\le\frac12\log m$, because $1/a\le\frac12\int_{a-2}^a dx/x$ and the integrals telescope. Hence $C_m(1)\le\frac{2m}\pi\bigl(1+\frac12\log m\bigr)=\frac m\pi(\log m+2)$.
\end{proof}

\begin{proposition}\label{prop:upper}
For every odd prime $p$,
\[
|S(p)|\;\le\;\frac2p\,C_p(1)^2\;\le\;\frac2{\pi^2}\,p\,(\log p+2)^2.
\]
\end{proposition}

\begin{proof}
By~\eqref{eq:dual}, Lemma~\ref{lem:sign-form} and the positivity of $\cot(\pi k/2p)$ for $0<k<p$,
\[
|S(p)|\le\frac2p\sum_{\substack{0<k<p\\ k\text{ odd}}}\cot\frac{\pi k}{2p}\,|C_p(k)|\le\frac2p\,C_p(1)\sum_{\substack{0<k<p\\ k\text{ odd}}}\cot\frac{\pi k}{2p}=\frac2p\,C_p(1)^2,
\]
since $C_p(1)=\sum_{k\text{ odd}}\cot\frac{\pi k}{2p}$. This is the first inequality. As $p\ge3$, Lemma~\ref{lem:Cm1-elementary} gives $C_p(1)\le\frac p\pi(\log p+2)$, and hence $\frac2pC_p(1)^2\le\frac2{\pi^2}p(\log p+2)^2$.
\end{proof}

\begin{corollary}
For all $C_2>0$ and $\varepsilon>0$ we have $S(p)<C_2p^{1+\varepsilon}$ for all sufficiently large primes $p$.
\end{corollary}

\begin{proof}
By Proposition~\ref{prop:upper}, $S(p)/p^{1+\varepsilon}\le\frac2{\pi^2}(\log p+2)^2p^{-\varepsilon}\to0$.
\end{proof}

\section{A reciprocity law for half-range cotangent sums}\label{sec:reciprocity}

The sums $C_p(k)$ in~\eqref{eq:dual} have length about $p/2$, like the original problem. The reciprocity law of this section (Lemma~\ref{lem:reciprocity}), proved by integrating a product of a tangent and a cotangent around a rectangle, writes $C_p(k)$ as $\frac pkF_k(p)+I(k/p)$, where $F_k(p)$ is a tangent sum of length $(k-1)/2$; Lemma~\ref{lem:I-asymptotics} gives the asymptotics of $I$, and Corollary~\ref{cor:Cm1} the resulting bounds for $C_m(1)$, $F_k(r)$ and $I(\lambda)$. Section~\ref{sec:decomposition} uses these results and the dual formula: the terms $I(k/p)$ produce the main term of $S(p)$, and the tangent sums produce its fluctuations.

\begin{lemma}[Reciprocity]\label{lem:reciprocity}
Let $m$ and $k$ be coprime odd positive integers. Then
\[
C_m(k)=\frac mk\,F_k(m)+I\Bigl(\frac km\Bigr).
\]
\end{lemma}

\begin{proof}
\emph{Idea.} The function $f(z)=-\frac\pi2\tan\frac{\pi z}2\cot\frac{\pi kz}{2m}$ has residues $\cot\frac{\pi ka}{2m}$ at the odd integers $a$ and $-\frac mk\tan\frac{\pi mn}{k}$ at the points $2mn/k$ in the strip $0<\operatorname{Re}z<m$. A contour integral over the boundary of the strip, a rectangle with height tending to infinity, turns the sum of the residues into $I(k/m)$.

\emph{Step 1: the residues in the strip.} The poles of $f$ with $0<\operatorname{Re}z<m$ are real and of two kinds.
\begin{enumerate}[label=(\roman*),leftmargin=2em]
\item The odd integers $a\in(0,m)$. There $\tan(\pi z/2)$ has a simple pole with residue $-2/\pi$, and $\cot(\pi ka/2m)$ is finite because $ka$ is odd. So the residue of $f$ is $\cot(\pi ka/2m)$.
\item The points $z_n=2mn/k$ with $1\le n\le(k-1)/2$, where $\cot(\pi kz/2m)$ has a simple pole with residue $2m/(\pi k)$. These points are not integers, because $k\nmid2mn$ (we have $(k,m)=1$ and $0<n<k$). Hence $\tan(\pi z_n/2)=\tan(\pi mn/k)$ is finite, and the residue of $f$ is $-\frac mk\tan(\pi mn/k)$.
\end{enumerate}
Since $C_m(k)=\sum_{a}\cot(\pi ka/2m)$ over odd $a\in(0,m)$ and $F_k(m)=\sum_n\tan(\pi mn/k)$, the sum of the residues in the strip is
\[
C_m(k)-\frac mkF_k(m).
\]

\emph{Step 2: the function on the boundary.} On the lines $\operatorname{Re}z=0$ and $\operatorname{Re}z=m$ the function $f$ is analytic. At $z=0$ a simple zero of $\tan(\pi z/2)$ cancels a simple pole of $\cot(\pi kz/2m)$. At $z=m$ a simple pole of $\tan(\pi z/2)$ ($m$ odd) cancels a simple zero of $\cot(\pi kz/2m)$ ($k$ odd). There are no other singularities on these lines. For real $y\ne0$,
\[
f(iy)=-\frac\pi2\tanh\frac{\pi y}2\coth\frac{\pi ky}{2m},\qquad f(m+iy)=-\frac\pi2\coth\frac{\pi y}2\tanh\frac{\pi ky}{2m}.
\]
Both functions of $y$ decay exponentially as $|y|\to\infty$ after subtracting their limit $-\pi/2$, and the differences $f(m+iy)-f(iy)$ are integrable over $\R$. As $|\operatorname{Im}z|\to\infty$ we have $\tan(\pi z/2)\to\pm i$ and $\cot(\pi kz/2m)\to\mp i$, uniformly in $\operatorname{Re}z$ and with exponentially small error, so $f(z)\to-\pi/2$.

\emph{Step 3: the contour integral.} Integrate $f$ counterclockwise over the boundary of the rectangle with vertices $-iT$, $m-iT$, $m+iT$, $iT$. As $T\to\infty$ the bottom side contributes $\to-\frac{\pi m}2$ and the top side, traversed from right to left, contributes $\to+\frac{\pi m}2$. The two vertical sides contribute $i\int_{-T}^T\bigl(f(m+iy)-f(iy)\bigr)\,dy$. By the residue theorem and Step~2,
\[
C_m(k)-\frac mkF_k(m)=\frac1{2\pi i}\cdot i\cdot\frac\pi2\int_\R\Bigl(\tanh\frac{\pi y}2\coth\frac{\pi ky}{2m}-\coth\frac{\pi y}2\tanh\frac{\pi ky}{2m}\Bigr)\,dy.
\]

\emph{Step 4: the integral is $I(k/m)$.} For $u\ne0$ and $\lambda>0$ we have the identity
\begin{equation}\label{eq:tanh-identity}
\tanh u\coth\lambda u-\coth u\tanh\lambda u=\frac{2(\cosh2u-\cosh2\lambda u)}{\sinh2u\,\sinh2\lambda u}.
\end{equation}
Indeed, the left side equals $(\sinh^2u\cosh^2\lambda u-\cosh^2u\sinh^2\lambda u)/(\sinh u\cosh u\sinh\lambda u\cosh\lambda u)$. The numerator equals $\cosh^2u-\cosh^2\lambda u=\frac12(\cosh2u-\cosh2\lambda u)$, and the denominator equals $\frac14\sinh2u\sinh2\lambda u$. Put $\lambda=k/m$ and $u=\pi y/2$ in Step~3. The integrand is even in $y$, the prefactor is $\frac14$, and $2u=\pi y$. Hence the right-hand side of Step~3 equals
\[
\frac14\cdot2\int_0^\infty\frac{2(\cosh\pi y-\cosh\lambda\pi y)}{\sinh\pi y\,\sinh\lambda\pi y}\,dy=\frac1\pi\int_0^\infty\frac{\cosh v-\cosh\lambda v}{\sinh v\,\sinh\lambda v}\,dv=I(\lambda),
\]
with $v=\pi y$.
\end{proof}

\begin{lemma}\label{lem:I-asymptotics}
For $\lambda>0$ write $\pi\lambda I(\lambda)=\log\frac1\lambda+\kz-\rho(\lambda)$. Then $0\le\rho(\lambda)\le\frac{7\pi^2}{72}\lambda^2$. Moreover $I(1/\lambda)=-\lambda I(\lambda)$.
\end{lemma}

\begin{proof}
\emph{Idea.} After the substitution $v=2u$, the integrand of $\pi I(\lambda)$ splits into a main part, whose integral is a Frullani integral and gives $\lambda^{-1}(\log\frac1\lambda+\kz)$, and two nonnegative remainders of size $O(1)$, which produce $\rho$.

\emph{Step 1: a decomposition of the integrand.} By~\eqref{eq:tanh-identity} and the substitution $v=2u$,
\[
\pi I(\lambda)=\int_0^\infty\bigl(\tanh u\coth\lambda u-\coth u\tanh\lambda u\bigr)\,du.
\]
Write $\coth x=\frac1x+\varkappa(x)$, where $0\le\varkappa(x)\le x/3$ for $x>0$ (from $\varkappa(x)=\sum_{n\ge1}\frac{2x}{x^2+\pi^2n^2}$). Put $\phi(x)=x/\sinh x$ and $\psi(x)=2/(e^x+1)=1-\tanh(x/2)$. Then
\begin{multline*}
\tanh u\coth\lambda u-\coth u\tanh\lambda u
=\frac{\phi(2\lambda u)-\psi(2u)}{\lambda u}\\
-(1-\tanh u)\,\varkappa(\lambda u)-(\coth u-1)\tanh\lambda u.
\end{multline*}
To check this, write $\tanh u\coth\lambda u=\frac{\tanh u}{\lambda u}+\varkappa(\lambda u)-(1-\tanh u)\varkappa(\lambda u)$ and $\coth u\tanh\lambda u=\tanh\lambda u+(\coth u-1)\tanh\lambda u$. Then use $\tanh u=1-\psi(2u)$ and $\varkappa(\lambda u)-\tanh\lambda u=\coth\lambda u-\tanh\lambda u-\frac1{\lambda u}=\frac2{\sinh2\lambda u}-\frac1{\lambda u}=\frac{\phi(2\lambda u)-1}{\lambda u}$. All three terms are integrable on $(0,\infty)$, and the last two are nonnegative.

\emph{Step 2: the Frullani integral.} Since $\phi$ is continuous with $\phi(0)=1$ and $\phi(\infty)=0$, Frullani's theorem gives
\[
\begin{aligned}
\int_0^\infty\bigl(\phi(2\lambda u)-\psi(2u)\bigr)\frac{du}u&=\int_0^\infty\bigl(\phi(2\lambda u)-\phi(2u)\bigr)\frac{du}u+\int_0^\infty\bigl(\phi(2u)-\psi(2u)\bigr)\frac{du}u\\
&=\log\frac1\lambda+\int_0^\infty\bigl(\phi(x)-\psi(x)\bigr)\frac{dx}x.
\end{aligned}
\]

\emph{Step 3: the value of the constant.} For $\operatorname{Re}s>0$, expanding the integrands in powers of $e^{-x}$ gives the Mellin transforms
\begin{align*}
\int_0^\infty\phi(x)x^{s-1}dx&=2\Gamma(s+1)(1-2^{-s-1})\zeta(s+1)=\frac1s+\log2+O(s),\\
\int_0^\infty\psi(x)x^{s-1}dx&=2\Gamma(s)(1-2^{1-s})\zeta(s)=\frac1s+\log\frac\pi2-\gamma+O(s).
\end{align*}
For the second expansion we used $\zeta(0)=-\frac12$ and $\zeta'(0)=-\frac12\log2\pi$. Since $\phi-\psi=O(x)$ near $0$ and decays exponentially at $\infty$, letting $s\to0^+$ in the difference gives
\[
\int_0^\infty(\phi-\psi)\,\frac{dx}x=\log2-\log\frac\pi2+\gamma=\kz.
\]

\emph{Step 4: the remainders.} Steps 1--3 give $\pi\lambda I(\lambda)=\log\frac1\lambda+\kz-\rho(\lambda)$ with $\rho(\lambda)=\lambda(J_1+J_2)$, where
\begin{gather*}
J_1=\int_0^\infty(1-\tanh u)\varkappa(\lambda u)\,du,\qquad J_2=\int_0^\infty(\coth u-1)\tanh(\lambda u)\,du.
\end{gather*}
Both are nonnegative. Using $\varkappa(x)\le x/3$, $\tanh x\le x$, $1-\tanh u=2\sum_{n\ge1}(-1)^{n-1}e^{-2nu}$ and $\coth u-1=2\sum_{n\ge1}e^{-2nu}$, we get
\begin{gather*}
J_1\le\frac\lambda3\int_0^\infty u(1-\tanh u)\,du=\frac{\pi^2\lambda}{72},\qquad
J_2\le\lambda\int_0^\infty u(\coth u-1)\,du=\frac{\pi^2\lambda}{12}.
\end{gather*}
Thus $0\le\rho(\lambda)\le(\frac{\pi^2}{72}+\frac{\pi^2}{12})\lambda^2=\frac{7\pi^2}{72}\lambda^2$.

\emph{Step 5: the symmetry.} The substitution $v\mapsto\lambda v$ in~\eqref{eq:I} gives $I(1/\lambda)=-\lambda I(\lambda)$.
\end{proof}

\begin{corollary}\label{cor:Cm1}
\textup{(a)} For odd $m\ge1$ we have $F_m(1)=C_m(1)=I(1/m)$.

\textup{(b)} For odd $k\ge3$ and every integer $r$, $|F_k(r)|\le F_k(1)\le\frac k\pi(\log k+2)$.

\textup{(c)} For $0<\lambda\le1$, $|I(\lambda)|\le\frac1{\pi\lambda}\bigl(\log\frac1\lambda+2\bigr)$.
\end{corollary}

\begin{proof}
(a) The substitution $a=m-2t$ gives
\[
F_m(1)=\sum_{1\le t<m/2}\tan\frac{\pi t}m=\sum_{\substack{0<a<m\\ a\text{ odd}}}\cot\frac{\pi a}{2m}=C_m(1),
\]
and Lemma~\ref{lem:reciprocity} with $k=1$ gives $C_m(1)=mF_1(m)+I(1/m)=I(1/m)$, since $F_1=0$.

(b) \emph{Idea.} If $(r,k)=1$, then $F_k(r)$ is a signed rearrangement of $F_k(1)$. In general $F_k(r)=F_{k'}(r')$ with $k'=k/(r,k)$ by periodicity and oddness, and $F_{k'}(1)\le F_k(1)$ by (a) and Lemma~\ref{lem:I-asymptotics}. The upper bound for $F_k(1)$ is (a) together with Lemma~\ref{lem:Cm1-elementary}.

\emph{Step 1: the upper bound.} By (a) and Lemma~\ref{lem:Cm1-elementary}, $F_k(1)=C_k(1)\le\frac k\pi(\log k+2)$ for odd $k\ge3$.

\emph{Step 2: coprime $r$.} Let $(r,k)=1$. The classes $\pm rt\bmod k$, $1\le t\le(k-1)/2$, are distinct, so $t\mapsto x(t)$, where $x(t)\in[1,\frac{k-1}2]$ and $x(t)\equiv\pm rt\pmod k$, is a bijection of $\{1,\dots,\frac{k-1}2\}$. Since $\tan(\pi(k-x)/k)=-\tan(\pi x/k)$, we get $F_k(r)=\sum_{x=1}^{(k-1)/2}\pm\tan(\pi x/k)$. All these tangents are positive, so $|F_k(r)|\le F_k(1)$.

\emph{Step 3: reduction of general $r$.} Let $g=(r,k)>1$, and put $k'=k/g$ and $r'=r/g$, so that $(r',k')=1$, $k'$ is odd and $g\ge3$. The function $t\mapsto\tan(\pi r't/k')$ has period $k'$ and is odd, so its sum over a complete period vanishes. Since $(k-1)/2=k'\frac{g-1}2+\frac{k'-1}2$, the sum over $1\le t\le(k-1)/2$ consists of $\frac{g-1}2$ complete periods and the sum over $1\le t\le(k'-1)/2$. Hence $F_k(r)=F_{k'}(r')$, and by Step~2, $|F_k(r)|\le F_{k'}(1)$. (If $k'=1$ this is $0$, and $F_k(1)\ge0$ is a sum of positive tangents, so $|F_k(r)|\le F_k(1)$.)

\emph{Step 4: $F_{k'}(1)\le F_k(1)$ for $k'\ge3$.} Here $k=gk'\ge3k'\ge9$. By (a) and Lemma~\ref{lem:I-asymptotics} with $\lambda=1/m$, for odd $m\ge3$
\[
\pi F_m(1)=\pi I(1/m)=m\bigl(\log m+\kz-\rho(1/m)\bigr),\qquad 0\le\rho(1/m)\le\frac{7\pi^2}{72m^2}.
\]
Hence $\pi F_{k'}(1)\le k'(\log k'+\kz)$ and $\pi F_k(1)\ge k(\log k+\kz)-\frac{7\pi^2}{72k}$. The function $x\mapsto x(\log x+\kz)$ is increasing for $x\ge1$ (its derivative is $\log x+\kz+1>0$), so $k\ge3k'$ gives $k(\log k+\kz)\ge3k'(\log k'+\log3+\kz)$. Since $\log k'+\kz>0$,
\[
\pi F_k(1)-\pi F_{k'}(1)\ge3k'(\log k'+\kz)+3k'\log3-k'(\log k'+\kz)-\frac{7\pi^2}{72k}\ge3k'\log3-\frac{7\pi^2}{72k}>0,
\]
because $3k'\log 3\ge9\log3>9$ and $\frac{7\pi^2}{72k}<1$. Together with Step~3 this gives $|F_k(r)|\le F_k(1)$ for all $r$, and with Step~1 the claim.

(c) By Lemma~\ref{lem:I-asymptotics}, $\pi\lambda I(\lambda)=\log\frac1\lambda+\kz-\rho(\lambda)$ with $0\le\rho(\lambda)\le\frac{7\pi^2}{72}\lambda^2<1$ for $\lambda\le1$. As $\log\frac1\lambda\ge0$ and $0<\kz<1$, we get $-1<\pi\lambda I(\lambda)<\log\frac1\lambda+\kz<\log\frac1\lambda+2$, hence $|\pi\lambda I(\lambda)|\le\log\frac1\lambda+2$.
\end{proof}

\section{The exact decomposition}\label{sec:decomposition}

Inserting the reciprocity law (Lemma~\ref{lem:reciprocity}) into the dual formula~\eqref{eq:dual} gives an exact identity $S(p)=\Psi(p)+\Tpart(p)$, in which the smooth part $\Psi(p)$ is an explicit sum equal to $\frac p2\log p+c_0p+O(1)$ (the source of the main term) and $\Tpart(p)$ is a weighted sum of the tangent sums $F_k(p)$. Only the moduli $k<p^{3/4}(\log p)^4$ matter in $\Tpart(p)$: for larger $k$ the sums $C_p(k)$ add up, over $k$ in an interval, to incomplete Kloosterman sums, and Kloosterman's bound $2p^{3/4}$, proved in full below, shows that these $k$ contribute $O(p/\log p)$. The resulting formula~\eqref{eq:decomp-ii} reduces the asymptotics of $S(p)$ to the study of $\sum_{k<p^{3/4}(\log p)^4}F_k(p)/k^2$, which is carried out in Sections~\ref{sec:kloosterman}--\ref{sec:proof}.

\begin{theorem}\label{thm:decomp}
\textup{(i)} For every odd prime $p$ we have the identity $S(p)=\Psi(p)+\Tpart(p)$, where
\[
\Psi(p)=\frac2p\sum_{\substack{0<k<p\\ k\text{ odd}}}\cot\frac{\pi k}{2p}\,I\Bigl(\frac kp\Bigr),\qquad\Tpart(p)=2\sum_{\substack{3\le k<p\\ k\text{ odd}}}\cot\frac{\pi k}{2p}\,\frac{F_k(p)}k,
\]
and $\Psi(p)=\frac p2\log p+c_0p+O(1)$.

\textup{(ii)} For every odd prime $p$ and every real $K$ with $p^{3/4}(\log p)^4\le K\le p^{9/10}$,
\begin{equation}\label{eq:decomp-ii}
S(p)=\frac p2\log p+c_0p+\frac{4p}\pi R_K(p)+O\Bigl(\frac p{\log p}\Bigr),\qquad R_K(p)=\sum_{\substack{3\le k<K\\ k\text{ odd}}}\frac{F_k(p)}{k^2}.
\end{equation}
The same holds with $R_K(p)$ replaced by the sum of $F_k(p)/k^2$ over all odd $k$ with $3\le k<p$. All implied constants are absolute.
\end{theorem}

The range of $K$ in~(ii) is empty unless $(\log p)^4\le p^{3/20}$, which forces $\log p>129$, that is, $p>10^{56}$; so~(ii) is a statement about very large $p$, and the asymptotic content is the $O(p/\log p)$ error. In part~(i), $\Psi(p)$ depends on $p$ only through the numbers $k/p$, and in $\Tpart(p)$ the sum $F_k(p)$ depends only on the residue of $p$ modulo $k$. In part~(ii) the upper limit $p^{9/10}$ is a matter of convenience: for $k<K$ the proof replaces $\frac2k\cot\frac{\pi k}{2p}$ by $\frac{4p}{\pi k^2}$, at a cost $O(K^2\log K/p)$.

\begin{proof}[Proof of Theorem~\textup{\ref{thm:decomp}(i)}]
\emph{Idea.} The identity is Lemma~\ref{lem:reciprocity} inserted into~\eqref{eq:dual}. For the evaluation of $\Psi(p)$, Lemma~\ref{lem:I-asymptotics} replaces $I(k/p)$ by its leading terms. The resulting error is $O(1)$, and the sum over $k$ is completed to all odd $k$ and evaluated by the Dirichlet series $\sum_{k\text{ odd}}k^{-s}=(1-2^{-s})\zeta(s)$.

\emph{The identity.} By Lemma~\ref{lem:reciprocity} with $m=p$ we have $C_p(k)=\frac pkF_k(p)+I(k/p)$ for odd $k<p$. Inserting this into~\eqref{eq:dual} gives $S(p)=\Psi(p)+\Tpart(p)$; the term $k=1$ does not contribute to $\Tpart(p)$ because $F_1=0$.

\emph{Reduction to an explicit sum.} Since $\cot\frac{\pi k}{2p}=\frac{2p}{\pi k}h(\frac kp)$, Lemma~\ref{lem:I-asymptotics} gives
\[
\Psi(p)=\frac{4p}{\pi^2}\sum_{\substack{0<k<p\\ k\text{ odd}}}\frac{h(k/p)}{k^2}\Bigl(\log\frac pk+\kz-\rho\Bigl(\frac kp\Bigr)\Bigr).
\]
Put $A_k=\log\frac pk+\kz\ (>0)$. Since $0\le1-h(x)\le x^2$ for $0<x<1$, $0\le h\le1$, and $0\le\rho(x)\le\frac{7\pi^2}{72}x^2$, replacing $h(\frac kp)\bigl(A_k-\rho(\frac kp)\bigr)$ by $A_k$ changes it by $A_k(1-h)+h\rho$, which lies between $0$ and
\[
\Bigl(\frac kp\Bigr)^2\Bigl(A_k+\frac{7\pi^2}{72}\Bigr)\le\Bigl(\frac kp\Bigr)^2\Bigl(\log\frac pk+2\Bigr),
\]
because $\kz+\frac{7\pi^2}{72}<2$. So the total change in $\Psi(p)$ is at most
\[
\frac{4}{\pi^2p}\sum_{k<p}\Bigl(\log\frac pk+2\Bigr)\le\frac{12}{\pi^2},
\]
since $\sum_{k\le p}\log\frac pk=p\log p-\log p!\le p$. Next, extending the sum over $k$ to all odd $k\ge1$ costs at most
\[
\frac{4p}{\pi^2}\sum_{k\ge p}\Bigl(\log\frac kp+\kz\Bigr)k^{-2}=O(1).
\]
Hence
\[
\Psi(p)=\frac{4p}{\pi^2}\sum_{k\text{ odd}}\frac{\log p+\kz-\log k}{k^2}+O(1)=\frac{4p}{\pi^2}\Bigl(\frac{\pi^2}{8}(\log p+\kz)-\sum_{k\text{ odd}}\frac{\log k}{k^2}\Bigr)+O(1),
\]
where we used $\sum_{k\text{ odd}}k^{-2}=\pi^2/8$.

\emph{The constant.} Differentiating $\sum_{k\text{ odd}}k^{-s}=(1-2^{-s})\zeta(s)$ at $s=2$ gives
\[
\sum_{k\text{ odd}}\frac{\log k}{k^2}=-\frac{\log2}{4}\zeta(2)-\frac34\zeta'(2).
\]
Since $\zeta(2)=\pi^2/6$, therefore $\Psi(p)=\frac p2(\log p+\kz)+\frac{p\log2}6+\frac{3p}{\pi^2}\zeta'(2)+O(1)=\frac p2\log p+c_0p+O(1)$, with $c_0$ as in~\eqref{eq:c0}.
\end{proof}

\subsection*{Kloosterman's bound}

For an integer $b$ with $p\nmid b$ and an integer $\nu$ put
\[
\operatorname{Kl}(b,\nu;p)=\sum_{z=1}^{p-1}e\Bigl(\frac{bz^*+\nu z}{p}\Bigr),
\]
where $z^*$ is the inverse of $z$ modulo $p$. The bound $|\operatorname{Kl}(b,\nu;p)|\le2p^{3/4}$ is due to Kloosterman; it is weaker than Weil's bound $2\sqrt p$ for $p\nmid\nu$, and we prove it in full.

\begin{lemma}[Kloosterman's bound; cf.\ \cite{b36}]\label{lem:kloosterman}
Let $p$ be an odd prime, $p\nmid b$ and $\nu\in\Z$. Then
\[
|\operatorname{Kl}(b,\nu;p)|\le\bigl(3p^2(p-2)\bigr)^{1/4}\le2p^{3/4}.
\]
\end{lemma}

\begin{proof}
\emph{Idea.} The sum depends only on $b\nu$ modulo $p$. The fourth moment of $\operatorname{Kl}(a,c;p)$ over all pairs $(a,c)$ modulo $p$ equals $p^2$ times the number of solutions of a pair of equations in $(\mathbb F_p^\times)^4$, and these solutions can be listed. Since each value $\operatorname{Kl}(1,\nu;p)$ is taken at $p-1$ pairs $(a,c)$, the fourth moment bounds it.

Throughout, $\operatorname{Kl}(a,c)=\operatorname{Kl}(a,c;p)$ is defined by the same formula for all $a,c\in\mathbb F_p$ (including $a=0$), and $x^*$ denotes the inverse in $\mathbb F_p^\times$.

\emph{Step 1: reduction to $\operatorname{Kl}(1,\cdot)$.} Let $b\in\mathbb F_p^\times$. The substitution $z=bw$ permutes $\mathbb F_p^\times$, and $bz^*=bb^*w^*=w^*$, $cz=cbw$. Hence $\operatorname{Kl}(b,c)=\operatorname{Kl}(1,bc)$ for every $c\in\mathbb F_p$. Therefore $\operatorname{Kl}(b,\nu)$ depends only on $b\nu\bmod p$, and it suffices to prove $|\operatorname{Kl}(1,a)|^4\le3p^2(p-2)$ for all $a\in\mathbb F_p$. (For $a=0$ this is the elementary evaluation $\operatorname{Kl}(1,0)=\sum_{z\ne0}e(z^*/p)=-1$.)

\emph{Step 2: the fourth moment.} For $x=(x_1,x_2,x_3,x_4)\in(\mathbb F_p^\times)^4$ put
\[
A(x)=x_1^*+x_2^*-x_3^*-x_4^*,\qquad B(x)=x_1+x_2-x_3-x_4.
\]
Writing $|\operatorname{Kl}(a,c)|^4=\operatorname{Kl}(a,c)^2\,\overline{\operatorname{Kl}(a,c)}^{\,2}$ and expanding,
\[
|\operatorname{Kl}(a,c)|^4=\sum_{x\in(\mathbb F_p^\times)^4}e\Bigl(\frac{aA(x)+cB(x)}p\Bigr).
\]
Summing over $a,c\in\mathbb F_p$ and using $\sum_{a\in\mathbb F_p}e(ar/p)=p\cdot\mathbf 1[r=0]$ gives
\[
\sum_{a,c\in\mathbb F_p}|\operatorname{Kl}(a,c)|^4=p^2N,\qquad N=\#\{x\in(\mathbb F_p^\times)^4:\ A(x)=0,\ B(x)=0\}.
\]

\emph{Step 3: counting the solutions.} We show $N=3(p-1)(p-2)$. The conditions say $x_1+x_2=x_3+x_4=:s$ and $x_1^*+x_2^*=x_3^*+x_4^*$. Since $x_1^*+x_2^*=s\,(x_1x_2)^*$ and $x_3^*+x_4^*=s\,(x_3x_4)^*$, the second condition reads $s\,(x_1x_2)^*=s\,(x_3x_4)^*$.

\emph{Case $s=0$.} The conditions reduce to $x_2=-x_1$ and $x_4=-x_3$, with $x_1,x_3\in\mathbb F_p^\times$ arbitrary. This gives $(p-1)^2$ solutions.

\emph{Case $s\ne0$.} The second condition is equivalent to $x_1x_2=x_3x_4$. Then $x_3,x_4$ are the two roots of $X^2-sX+x_1x_2$, which has $x_1,x_2$ as roots; so $(x_3,x_4)\in\{(x_1,x_2),(x_2,x_1)\}$. Conversely both choices are solutions. The pairs $(x_1,x_2)\in(\mathbb F_p^\times)^2$ with $x_1+x_2\ne0$ number $(p-1)^2-(p-1)=(p-1)(p-2)$. For each of them there are two choices of $(x_3,x_4)$, which coincide exactly when $x_1=x_2$; and $x_1=x_2$ with $s=2x_1\ne0$ holds for $p-1$ pairs (as $p$ is odd). So this case gives $2(p-1)(p-2)-(p-1)=(p-1)(2p-5)$ solutions.

The cases are disjoint, so $N=(p-1)^2+(p-1)(2p-5)=(p-1)(3p-6)=3(p-1)(p-2)$.

\emph{Step 4: conclusion.} Fix $a\in\mathbb F_p$. For each $b\in\mathbb F_p^\times$ the pair $(b,b^*a)$ satisfies $\operatorname{Kl}(b,b^*a)=\operatorname{Kl}(1,bb^*a)=\operatorname{Kl}(1,a)$ by Step~1, and the $p-1$ pairs are distinct. All terms in the fourth moment are nonnegative, so
\[
(p-1)\,|\operatorname{Kl}(1,a)|^4\le\sum_{a',c\in\mathbb F_p}|\operatorname{Kl}(a',c)|^4=3p^2(p-1)(p-2).
\]
Hence $|\operatorname{Kl}(1,a;p)|^4\le3p^2(p-2)<3p^3$, so $|\operatorname{Kl}(1,a;p)|<3^{1/4}p^{3/4}<2p^{3/4}$. By Step~1 this proves the lemma for all $b$ with $p\nmid b$ and all $\nu$.
\end{proof}

\begin{lemma}\label{lem:incomplete}
Let $p$ be an odd prime, $c$ an integer with $p\nmid c$, and $J\subset[1,p)$ an interval. Then
\[
\Big|\sum_{\substack{k\in J\\ k\text{ odd}}}e\Bigl(\frac{c\overline{k}}{2p}\Bigr)\Big|\le2p^{3/4}(\log p+2).
\]
\end{lemma}

\begin{proof}
\emph{Step 1: reduction modulo $p$.} Let $2^*=(p+1)/2$, the inverse of $2$ modulo $p$. Since $\frac1{2p}=\frac{2^*}{p}-\frac12$, we have $e(cx/2p)=(-1)^ce(c\,2^*x/p)$ for odd $x$. Apply this to the odd residue $x=\overline{k}$, which is $\equiv k^*\pmod p$. This gives $e(c\overline{k}/2p)=(-1)^ce(c\,2^*k^*/p)$. For $k=2j+1$ and $y=j+2^*$ we have $2y=2j+p+1\equiv k\pmod p$, so $k^*\equiv2^*y^*$ and $y\not\equiv0$. Hence the sum in question equals $\pm\sum_{y\in J'}e(by^*/p)$ with $b=c\,(2^*)^2\not\equiv0\pmod p$ and $J'=\{j+2^*:\ 2j+1\in J\}$. This is a set of at most $(p-1)/2$ consecutive integers, none of them divisible by $p$.

\emph{Step 2: completion.} Since $J'$ consists of at most $p$ consecutive integers, none divisible by $p$, orthogonality gives
\[
\sum_{y\in J'}e\Bigl(\frac{by^*}{p}\Bigr)=\frac1p\sum_{s\bmod p}\Bigl(\sum_{y\in J'}e\Bigl(\frac{sy}{p}\Bigr)\Bigr)
\operatorname{Kl}(b,-s;p).
\]

\emph{Step 3: the bound.} Here $|\operatorname{Kl}(b,0;p)|=1$, and $|\operatorname{Kl}(b,-s;p)|\le2p^{3/4}$ for $p\nmid s$ by Lemma~\ref{lem:kloosterman}. Since $|\sum_{y\in J'}e(sy/p)|\le1/(2\|s/p\|)$ and $\sum_{s=1}^{p-1}1/\|s/p\|=2p\sum_{s\le(p-1)/2}1/s\le2p(1+\log p)$, the sum is at most
\[
\frac{|J'|}{p}+\frac{2p^{3/4}}{p}\cdot\frac12\sum_{s=1}^{p-1}\frac1{\|s/p\|}\le1+2p^{3/4}(1+\log p)\le2p^{3/4}(\log p+2),
\]
where the last step uses $1\le2p^{3/4}$.
\end{proof}

\begin{lemma}\label{lem:Cp-interval}
For every odd prime $p$ and every interval $J\subset[1,p)$,
\[
\Big|\sum_{\substack{k\in J\\ k\text{ odd}}}C_p(k)\Big|\le\frac4{\pi^2}\,p^{7/4}(\log p+2)^3.
\]
\end{lemma}

\begin{proof}
Let $\tilde\sigma$ be the function on $\Z/2p\Z$ that equals $\sigma$ outside $\{0,p\}$ and vanishes at $0$ and $p$. Then $\tilde\sigma=f-\delta_0+\delta_p$ with $f$ as in the proof of Lemma~\ref{lem:dual}. By~\eqref{eq:fourier-coeff}, the Fourier coefficients of $\tilde\sigma$ vanish for even frequencies, and for odd $s$ they equal $\hat f(s)-\frac1p=-\frac ip\cot\frac{\pi s}{2p}$. In particular the coefficient at $s=p$ vanishes, and the sum of the absolute values of all coefficients is $\frac2pC_p(1)$. By~\eqref{eq:sign-form},
\[
\sum_{\substack{k\in J\\ k\text{ odd}}}C_p(k)=\sum_{\substack{0<n<p\\ n\text{ odd}}}\cot\frac{\pi n}{2p}
\sum_{\substack{0<s<2p,\ s\ne p\\ s\text{ odd}}}\hat{\tilde\sigma}(s)
\sum_{\substack{k\in J\\ k\text{ odd}}}e\Bigl(\frac{sn\overline{k}}{2p}\Bigr).
\]
Since $0<s<2p$, $s\ne p$ and $0<n<p$, we have $p\nmid sn$, and Lemma~\ref{lem:incomplete} bounds the inner sum by $2p^{3/4}(\log p+2)$. Therefore the left-hand side is at most
\[
C_p(1)\cdot\frac2pC_p(1)\cdot2p^{3/4}(\log p+2),
\]
and Lemma~\ref{lem:Cm1-elementary}, $C_p(1)\le\frac p\pi(\log p+2)$, gives $\frac{2}{p}\cdot\frac{p^2}{\pi^2}(\log p+2)^2\cdot2p^{3/4}(\log p+2)=\frac4{\pi^2}p^{7/4}(\log p+2)^3$.
\end{proof}

\begin{proof}[Proof of Theorem~\textup{\ref{thm:decomp}(ii)}]
\emph{Idea.} Split $\Tpart(p)$ at $K$. For $k\ge K$, write $F_k(p)$ through $C_p(k)$ by Lemma~\ref{lem:reciprocity}; the $I(k/p)$-terms are negligible, and the $C_p(k)$ are summed over dyadic intervals with Lemma~\ref{lem:Cp-interval}. For $k<K$ the weight $\frac2k\cot\frac{\pi k}{2p}$ is replaced by $\frac{4p}{\pi k^2}$, using the trivial bound for $F_k(p)$.

Let $K$ be in the range $p^{3/4}(\log p)^4\le K\le p^{9/10}$ (nonempty only for large $p$, as noted above) and write $\Tpart(p)=\Tpart_{<K}(p)+\Tpart_{\ge K}(p)$, where $\Tpart_{<K}$ and $\Tpart_{\ge K}$ are the parts of the sum with $k<K$ and $k\ge K$. Write $L=\log p$.

\emph{Step 1: large $k$, the terms $I(k/p)$.} By Lemma~\ref{lem:reciprocity}, for odd $k<p$,
\[
2\cot\frac{\pi k}{2p}\,\frac{F_k(p)}{k}=\frac2p\cot\frac{\pi k}{2p}\Bigl(C_p(k)-I\Bigl(\frac kp\Bigr)\Bigr).
\]
By Corollary~\ref{cor:Cm1}(c) and $\cot\frac{\pi k}{2p}\le\frac{2p}{\pi k}$, the terms $I(k/p)$ contribute at most
\[
\frac2p\sum_{K\le k<p}\frac{2p}{\pi k}\cdot\frac{p}{\pi k}\Bigl(\log\frac pk+2\Bigr)\le\frac{4p}{\pi^2}(L+2)
\sum_{k\ge K}\frac1{k^2}\ll\frac{pL}{K}\le\frac{p^{1/4}}{L^3}\ll\frac p{L}.
\]

\emph{Step 2: large $k$, the terms $C_p(k)$.} Cover $[K,p)$ by the intervals $B_j=[2^jK,2^{j+1}K)\cap[K,p)$, $j\ge0$. On $B_j$ the weight $\cot(\pi k/2p)$ is positive and decreasing, and it is at most $2p/(\pi2^jK)$, since $\cot x\le1/x$. Abel summation and Lemma~\ref{lem:Cp-interval}, applied to the subintervals of $B_j$, give
\[
\Big|\sum_{k\in B_j}\cot\frac{\pi k}{2p}\,C_p(k)\Big|\le\frac{2p}{\pi2^jK}\cdot\frac4{\pi^2}\,p^{7/4}(L+2)^3.
\]
Summing over $j$ (with $\sum_{j\ge0}2^{-j}=2$) and multiplying by $\frac2p$, we obtain at most
\[
\frac{32}{\pi^3}\,\frac{p^{7/4}(L+2)^3}{K}\ll\frac{p^{7/4}L^3}{p^{3/4}L^4}=\frac pL.
\]
Together with Step~1 this gives $\Tpart_{\ge K}(p)\ll p/L$.

\emph{Step 3: small $k$.} For $3\le k<K$ we have $2\cot\frac{\pi k}{2p}\frac1k=\frac{4p}{\pi k^2}h(\frac kp)$ and $0\le1-h(\frac kp)\le\frac{k^2}{p^2}$. By Corollary~\ref{cor:Cm1}(b), $|F_k(p)|\le\frac k\pi(\log k+2)$, so
\[
\bigg|\Tpart_{<K}(p)-\frac{4p}{\pi}R_K(p)\bigg|\le\frac4{\pi p}\sum_{\substack{k<K\\ k\text{ odd}}}|F_k(p)|
\le\frac4{\pi^2p}\sum_{\substack{k<K\\ k\text{ odd}}}k(\log k+2)\ll\frac{K^2\log K}{p}\ll p^{4/5}L.
\]
Here we used $K\le p^{9/10}$, so that $K^2\le p^{9/5}$ and $\log K\le L$. Combining Steps 1--3 with part~(i), and using $p^{4/5}L\ll p/L$, we obtain~\eqref{eq:decomp-ii}.

\emph{Step 4: the sum over all $k<p$.} Take $K=p^{3/4}L^4$, which lies in the range of~\eqref{eq:decomp-ii} (for the $p$ for which that range is nonempty). By~\eqref{eq:decomp-ii} it suffices to show that
\[
\frac{4p}{\pi}\sum_{\substack{K\le k<p\\ k\text{ odd}}}\frac{F_k(p)}{k^2}\ll\frac p{L}.
\]
By Lemma~\ref{lem:reciprocity}, this sum equals $\sum_{K\le k<p}\frac4{\pi k}\bigl(C_p(k)-I(\frac kp)\bigr)$. Steps~1 and~2 apply verbatim with the positive decreasing weight $\frac4{\pi k}$ in place of $\frac2p\cot\frac{\pi k}{2p}$: the bound $\frac2p\cot\frac{\pi k}{2p}\le\frac4{\pi k}$ was the only property of the weight used in Step~1, and in Step~2 the weight on $B_j$ is at most $\frac4{\pi2^jK}=\frac2p\cdot\frac{2p}{\pi2^jK}$, as before. This proves the last assertion of part~(ii).
\end{proof}

\begin{remark}\label{rem:weil}
With Weil's bound $|\operatorname{Kl}(b,\nu;p)|\le2\sqrt p$ for $p\nmid\nu$ in place of Lemma~\ref{lem:kloosterman}, the same proof gives Theorem~\ref{thm:decomp}(ii) for the cutoff $K\ge\sqrt p(\log p)^4$. This is not needed here, because the cutoff $p^{3/4}(\log p)^4$ is sufficient for the rest of the argument; Theorem~\ref{thm:decomp} uses no deep input. Weil's bound does enter the paper, however, through the bound for short Kloosterman sums used in Section~\ref{sec:kloosterman}: for sums of length $N\ge p^{4/7}$ the proof of that bound rests on it (see the proof of Lemma~\ref{lem:bg}).
\end{remark}

\section{Short Kloosterman sums}\label{sec:kloosterman}

For a prime $p$, an integer $b\not\equiv0\pmod p$ and $N<p$, the sums $\sum_{n\le N}e(bn^*/p)$, where $n^*$ is the inverse of $n$ modulo $p$, are the \emph{short Kloosterman sums}; by completion, Weil's bound for complete Kloosterman sums gives $\ll\sqrt p\log p$ for them, which is trivial for $N<\sqrt p$, and nontrivial bounds for smaller $N$ are known only when $\log N$ is at least about $(\log p)^{2/3}$, and then they save only a power of $\log p$ (see Garaev and Shparlinski~\cite[\S1.1]{b12}, Bourgain and Garaev~\cite{b8,b9} and Korolev~\cite{b17}; but see the remark on~\cite{b17} in~\cite[\S1.1]{b12}). The bound of Bourgain and Garaev~\cite{b8} is the only external input of this paper, and it is used here for $\log N\ge\frac12(\log p)^{9/10}$, far inside the range covered by its proof (Lemma~\ref{lem:bg}). This section defines the quantity $\mathcal{M}_K$, which measures cancellation in sums of $e(c\overline{k}/2p)$ over odd $k$ in a dyadic interval, and reduces it to the sums $\sum_{n\le N}e(bn^*/p)$ (Lemma~\ref{lem:reduction}); Corollary~\ref{cor:MK} then bounds $\mathcal{M}_K$ for $\log K\ge(\log p)^{9/10}$, and Lemma~\ref{lem:heads} in Section~\ref{sec:tangent} uses $\mathcal{M}_K$.

\subsection{The theorem of Bourgain and Garaev}
For a large prime $p$, Bourgain and Garaev~\cite[Theorem~16]{b8} state the following.
\begin{quote}
\emph{Theorem 16. The following bound holds:}
\[
\max_{(a,p)=1}\Bigl|\sum_{n\le N}e_p(an^*)\Bigr|\ll\frac{(\log\log p)^3\log p}{(\log N)^{3/2}}\,N,
\]
\emph{where the implied constant is absolute.}
\end{quote}
Here $e_p(x)=e(x/p)$. The statement has no hypothesis on $N$; it is nontrivial only when $(\log N)^{3/2}$ is larger than $(\log p)(\log\log p)^3$. We use it only in the range of $N$ that its proof covers, stated in the next lemma.

\begin{lemma}[Input from Bourgain--Garaev]\label{lem:bg}
There are absolute constants $C_{BG}$ and $p_0$ such that for every prime $p\ge p_0$, every $b\not\equiv0\pmod p$ and every $N<p$ with $\log N\ge\frac12(\log p)^{9/10}$,
\[
\Big|\sum_{1\le n\le N}e\Bigl(\frac{bn^*}{p}\Bigr)\Big|\le C_{BG}\,\frac{(\log\log p)^3\log p}{(\log N)^{3/2}}\,N.
\]
\end{lemma}

\begin{proof}
As stated in~\cite{b8}, Theorem~16 has no hypothesis on $N$, so the lemma is a special case of it. Since the theorem is used only in this range, we indicate why its proof covers the range, so that the reader can check this against~\cite{b8} (page numbers refer to arXiv v1, the only arXiv version).

The proof puts $\varepsilon=\log N/\log p$. For $\varepsilon\ge\frac47$ the bound follows from Weil's bound for incomplete Kloosterman sums (p.~54), since then $\sqrt p\log p\ll Np^{-1/14}\log p$, which is smaller than the right-hand side of the lemma. This case does occur: Section~\ref{sec:proof} applies the lemma with $N$ up to about $p^{3/4}(\log p)^4$. So Weil's bound enters the paper through Lemma~\ref{lem:bg}, although Kloosterman's weaker bound suffices in Section~\ref{sec:decomposition}. For $\varepsilon<\frac47$ the proof takes $\alpha=1/\log\log p$ and, at the end (p.~60),
\[
\beta=\frac{C\log\log p}{\varepsilon^{3/2}(\log p)^{1/2}}
\]
with a sufficiently large absolute constant $C$. On p.~55 it works with $\beta<(\log\log p)^{-3}$. This holds if and only if
\begin{align*}
\varepsilon^{3/2}>\frac{C(\log\log p)^4}{(\log p)^{1/2}}
&\iff\varepsilon>\frac{C^{2/3}(\log\log p)^{8/3}}{(\log p)^{1/3}}\\
&\iff\log N>C^{2/3}(\log p)^{2/3}(\log\log p)^{8/3}.
\end{align*}
In the range of the lemma, $\varepsilon\ge\frac12(\log p)^{-1/10}$, so
\[
C\log\log p\,(\log p)^{-1/2}\le\beta\le2^{3/2}C\log\log p\,(\log p)^{-7/20},
\]
and the other requirements of the proof hold for large $p$, for example:
\begin{itemize}[leftmargin=2em]
\item $\beta>2\log\log N/\log N$, because $\log N\ge\frac12(\log p)^{9/10}$;
\item $\alpha>\beta$, because $\beta\le(\log p)^{-1/3}$ for large $p$;
\item $\varepsilon\beta>(\log p)^{-1/2}$, because $\varepsilon\beta=C\log\log p/(\varepsilon^{1/2}(\log p)^{1/2})$ and $\varepsilon\le1$.
\end{itemize}
The remaining conditions on pp.~55--59 ($\alpha<0.1$; $\beta\log\log p<0.1$; de Bruijn's condition in the two estimates for smooth numbers; $\varepsilon\alpha>(\log p)^{-1/2}$; and $(2k_j)^{k_j}<p^{1/10}$) also hold for large $p$ in this range.
Finally, $\frac12(\log p)^{9/10}>C^{2/3}(\log p)^{2/3}(\log\log p)^{8/3}$ for large $p$, so every $N$ in the lemma lies in the range covered by the proof.
\end{proof}

\subsection{Reduction and the bound for \texorpdfstring{$\mathcal{M}_K$}{M\_K}}
For $6\le K\le p/2$ let
\begin{equation}\label{eq:MK}
\mathcal{M}_K=\max\Bigl\{\Big|\sum_{\substack{k\in I\\ k\text{ odd}}}e\Bigl(\frac{c\overline{k}}{2p}\Bigr)\Big|:\ p\nmid c,\
I\subset[K,2K)\text{ an interval}\Bigr\}.
\end{equation}
Here $\overline{k}$ is the inverse of $k$ modulo $2p$, as in Section~\ref{sec:notation}; the sum is over odd $k$ with $K\le k<2K\le p$, so $k$ is prime to $2p$.

\begin{lemma}[Reduction to initial sums]\label{lem:reduction}
Let $6\le K\le p/2$. Then
\[
\mathcal{M}_K\le4\max\Bigl\{\Big|\sum_{1\le n\le N'}e\Bigl(\frac{bn^*}{p}\Bigr)\Big|:\ p\nmid b,\ \frac K2-1\le N'<2K\Bigr\},
\]
where $N'$ runs over integers.
\end{lemma}

\begin{proof}
Let $2^*=(p+1)/2$, the inverse of $2$ modulo $p$. Since $\frac1{2p}=\frac{2^*}{p}-\frac12$, we have
$e(cx/2p)=(-1)^ce(c\,2^*x/p)$ for odd $x$. Apply this to the odd residue $x=\overline{k}$, which satisfies $\overline{k}\equiv k^*\pmod p$:
\[
e\Bigl(\frac{c\overline{k}}{2p}\Bigr)=(-1)^ce\Bigl(\frac{ak^*}{p}\Bigr),\qquad a=c\,2^*.
\]
For an interval $I=[u,v]\cap\Z\subset[K,2K)$ we have, because $(2j)^*=2^*j^*$,
\[
\sum_{\substack{k\in I\\ k\text{ odd}}}e\Bigl(\frac{ak^*}{p}\Bigr)=\sum_{k\in I}e\Bigl(\frac{ak^*}{p}\Bigr)-\sum_{u/2\le j\le v/2}e\Bigl(\frac{a\,2^*j^*}{p}\Bigr).
\]
Each of the two sums on the right is the difference of two initial sums $\sum_{n\le N'}e(bn^*/p)$, namely
\[
\sum_{k\in I}=\Sigma_{a}(v)-\Sigma_{a}(u-1),\qquad \sum_{u/2\le j\le v/2}=\Sigma_{a2^*}(\lfloor v/2\rfloor)-\Sigma_{a2^*}(\lceil u/2\rceil-1),
\]
where $\Sigma_b(N')=\sum_{1\le n\le N'}e(bn^*/p)$. Here $b\in\{a,a\,2^*\}$ is prime to $p$ because $p\nmid c$, and
$N'\in\{v,\,u-1,\,\lfloor v/2\rfloor,\,\lceil u/2\rceil-1\}\subset[\frac K2-1,2K)$; all these $N'$ are below $p$ since $v<2K\le p$.
The four initial sums give the factor $4$.
\end{proof}

\begin{corollary}\label{cor:MK}
There is an absolute constant $p_1$ such that for every prime $p\ge p_1$ and every $K$ with $\log K\ge(\log p)^{9/10}$ and $K\le p/2$,
\[
\frac{\mathcal{M}_K}K\le16\sqrt2\,C_{BG}\,\frac{(\log\log p)^3\log p}{(\log K)^{3/2}},
\]
where $C_{BG}$ is the constant of Lemma~\ref{lem:bg}.
\end{corollary}

\begin{proof}
Take $p_1\ge p_0$ so large that $K\ge e^{(\log p)^{9/10}}\ge9$ and $\log3\le\frac12(\log p)^{9/10}$ for $p\ge p_1$. Then $K\ge9>6$, and Lemma~\ref{lem:reduction} applies. Let $N'$ be an integer with $\frac K2-1\le N'<2K$; then $N'<2K\le p$.

\emph{Step 1: $N'$ is in the range of Lemma~\ref{lem:bg}.} In the proof of Lemma~\ref{lem:reduction} the lengths $N'$ that occur satisfy $N'\le v<2K\le p$, so $N'<p$. As $K\ge6$, $N'\ge\frac K2-1\ge\frac K3$, hence
\[
\log N'\ge\log K-\log3\ge(\log p)^{9/10}-\tfrac12(\log p)^{9/10}=\tfrac12(\log p)^{9/10}.
\]

\emph{Step 2: $\log N'\ge\frac12\log K$.} As $K\ge9$, $N'\ge\frac K2-1\ge\sqrt K$ (the function $\frac K2-1-\sqrt K$ is increasing for $K\ge1$ and equals $\frac12$ at $K=9$).

\emph{Step 3: the constant.} By Lemma~\ref{lem:bg}, Step~2 and $N'\le2K$,
\[
\frac{1}{K}\Big|\sum_{n\le N'}e\Bigl(\frac{bn^*}{p}\Bigr)\Big|\le\frac{N'}{K}\,C_{BG}\frac{(\log\log p)^3\log p}{(\log N')^{3/2}}\le2\cdot2^{3/2}\,C_{BG}\frac{(\log\log p)^3\log p}{(\log K)^{3/2}}.
\]
Lemma~\ref{lem:reduction} multiplies this by $4$, which gives the constant $4\cdot2\cdot2^{3/2}=16\sqrt2$.
\end{proof}

Section~\ref{sec:proof} uses Corollary~\ref{cor:MK} for the moduli $\exp((\log p)^{9/10})\le k<p^{3/4}(\log p)^4$.

\section{Heads and tails of the tangent sums}\label{sec:tangent}

This section provides the tools for the terms of~\eqref{eq:decomp-ii} with $k<p^{3/4}(\log p)^4$: a formula for
$F_k(p)/k$ in terms of $\theta_k$ alone, bounds for its head on average over $k$, and bounds for its tail for each $k$
separately. Section~\ref{sec:proof} combines them.

\emph{The split.} Lemma~\ref{lem:sign-form-tangent} writes $\Phi_k=F_k(p)/k$ as a sum over odd $n<k$ of the weights $c_n(k)\le\frac2{\pi n}$ times the signs $\sq(n\theta_k)$.
For a parameter $N$ the \emph{head} $\Phi_k^{\le N}$ is the part with $n\le N$, and the \emph{tail} $\Phi_k^{>N}$ is the rest.
The head is bounded trivially by $\log N$ (eq.~\eqref{eq:head-trivial}), which suffices for small $k$, and for larger $k$ is bounded on average over $k$ in a dyadic interval by Lemma~\ref{lem:heads}.
Its proof expands $\sq$ in a Fourier series (Lemmas~\ref{lem:vaaler} and~\ref{lem:square-wave}) and uses the bound for $\mathcal{M}_K$ from Section~\ref{sec:kloosterman}.
The tail is bounded for each $k$ separately: Lemma~\ref{lem:tails} reduces it to the quantity $D_k$, which is small when $\theta_k$ is not close to a rational number with small denominator (Lemma~\ref{lem:generic}).
Lemma~\ref{lem:witnesses} converts the opposite case into a divisibility $k\mid pr-m$ with $|r|$ small, which is handled by Theorem~\ref{thm:divisors} in Section~\ref{sec:divisors}.

\begin{lemma}\label{lem:sign-form-tangent}
Let $k$ be odd with $3\le k<p$, let $u_k\in[1,k)$ satisfy $2pu_k\equiv1\pmod k$, and put $\theta_k=u_k/k$.
\begin{enumerate}[label=\textup{(\alph*)},leftmargin=*]
\item $F_k(p)=\sum_{n\text{ odd},\,0<n<k}\cot\frac{\pi n}{2k}\,\sigma(n\overline{k})$.
\item For odd $n$ with $0<n<k$ we have $\sigma(n\overline{k})=-\sgn(\frac12-\{n\theta_k\})$; hence
\[
F_k(p)=-\sum_{\substack{0<n<k\\ n\text{ odd}}}\cot\frac{\pi n}{2k}\,\sgn\Bigl(\frac12-\{n\theta_k\}\Bigr)
\]
depends only on $\theta_k$.
\item $\theta_k\equiv\frac1{2pk}-\frac{\overline{k}}{2p}\pmod 1$.
\end{enumerate}
\end{lemma}

\begin{proof}
The idea is to identify, for each odd $n$, the residue $t$ with $pt\equiv(k-n)/2\pmod k$ with the integer $j$ defined by $kw=n+2pj$, where $w\equiv n\overline{k}\pmod{2p}$.
Then the sign $\sigma(n\overline{k})$ is read off from the position of $t$ relative to $k/2$.

(a) As $t$ runs over $1,\dots,k-1$, $x=pt\bmod k$ runs over $1,\dots,k-1$; write $t_x$ for the corresponding $t$, and
let $\delta(t)=1$ for $t\le(k-1)/2$ and $\delta(t)=-1$ otherwise. Since $\tan(\pi pt/k)=\tan(\pi x/k)$ and $\sum_{x=1}^{k-1}\tan(\pi x/k)=0$, we get
\[
F_k(p)=\sum_{x=1}^{k-1}\tan\frac{\pi x}k\cdot\frac{1+\delta(t_x)}{2}=\frac12\sum_{x=1}^{k-1}\tan\frac{\pi x}k\,\delta(t_x).
\]
As $t_{k-x}=k-t_x$, the terms $x$ and $k-x$ coincide, and writing $x=(k-n)/2$ with $n$ odd, $0<n<k$, we obtain
$F_k(p)=\sum_n\cot\frac{\pi n}{2k}\,\delta(t_{(k-n)/2})$, where $t_{(k-n)/2}\equiv-nu_k\pmod k$ (multiply $pt\equiv(k-n)/2$ by $2u_k$).
On the other hand,
let $w\in(0,2p)$ with $w\equiv n\overline{k}\pmod{2p}$. Then $kw=n+2pj$ for an integer $j$, and $0<kw<2pk$,
$0<n<k<p$ give $0\le j\le k-1$; moreover $j\equiv-nu_k\pmod k$, and $j\ne0$ since $j=0$ would give $kw=n$, so $k\mid n$. Thus $j=t_{(k-n)/2}$.
Finally $\sigma(n\overline{k})=1$ if and only if $w<p$, i.e.\ $n+2pj<pk$, i.e.\ $j<\frac k2-\frac n{2p}$; as
$0<\frac n{2p}<\frac12$ and $k$ is odd, this means $j\le\frac{k-1}2$. Hence $\sigma(n\overline{k})=\delta(j)$, proving (a).

(b) With $j$ as above, $j=k-(nu_k\bmod k)=k(1-\{n\theta_k\})$, so $j\le\frac{k-1}2$ if and only if
$\{n\theta_k\}>\frac12$ (note $\{n\theta_k\}\in\frac1k\Z$ and $k$ is odd, so $\{n\theta_k\}\neq\frac12$).
Thus $\sigma(n\overline{k})=1$ if and only if $\sgn(\frac12-\{n\theta_k\})=-1$.

(c) Since $2pu_k+k\overline{k}\equiv1$ modulo $k$ and modulo $2p$, we have $2pu_k+k\overline{k}=1+2pkm$ with
$m\in\Z$; divide by $2pk$.
\end{proof}

Throughout the rest of the paper we write, for odd $k$ with $3\le k<p$,
\begin{gather*}
\sq(y)=\sgn\Bigl(\frac12{-}\{y\}\Bigr)\ \ (y\notin\tfrac12\Z),\qquad
\Phi_k=\frac{F_k(p)}{k},\\
c_n(k)=\frac1k\cot\frac{\pi n}{2k}=\frac2{\pi n}\,h\Bigl(\frac nk\Bigr)\ \ (0<n<k),
\end{gather*}
where $\sq$ is the function of Section~\ref{sec:notation}. For $p\nmid m$ we have $\sigma(m)=\sq(m/2p)$, because $\{m/2p\}<\frac12$ exactly when $m\bmod2p\in(0,p)$. By
Lemma~\ref{lem:sign-form-tangent},
\begin{equation}\label{eq:Phik-sum}
\Phi_k=\sum_{\substack{0<n<k\\ n\text{ odd}}}c_n(k)\,\sigma(n\overline{k}),\qquad
\sigma(n\overline{k})=\sq\Bigl(\frac{n\overline{k}}{2p}\Bigr)=-\sq(n\theta_k),
\end{equation}
and none of the arguments of $\sq$ lies in $\frac12\Z$ (if $n\theta_k\in\frac12\Z$ then $k\mid2nu_k$, so $k\mid n$, since $\gcd(2u_k,k)=1$; and $p\nmid n\overline{k}$). For
$N\ge1$ we split $\Phi_k=\Phi_k^{\le N}+\Phi_k^{>N}$, where $\Phi_k^{\le N}$ is the part of the sum in
\eqref{eq:Phik-sum} with $n\le N$; thus $\Phi_k^{>N}=0$ if $k\le N+1$. Since $0\le h\le1$, we have
$0<c_n(k)\le\frac2{\pi n}$, and since $\sum_{n\text{ odd},\,3\le n\le X}\frac1n\le\frac12\log X$,
\begin{equation}\label{eq:head-trivial}
|\Phi_k^{\le N}|\le\frac2\pi\Bigl(1+\frac12\log N\Bigr).
\end{equation}

\begin{lemma}[Vaaler \cite{b27}; see {\cite[Theorem~A.6]{b13}}]\label{lem:vaaler}
For every integer $H\ge1$ there are complex numbers $a_\nu$, $1\le|\nu|\le H$, with $|a_\nu|\le\frac1{2\pi|\nu|}$,
such that for all real $x$
\[
\Big|\dpar{x}-\sum_{1\le|\nu|\le H}a_\nu e(\nu x)\Big|\le\Delta_H(x):=\frac1{2H+2}\sum_{|\nu|\le H}
\Bigl(1-\frac{|\nu|}{H+1}\Bigr)e(\nu x).
\]
\end{lemma}

Note that $\Delta_H\ge0$ (it is a multiple of the Fej\'er kernel).

\begin{lemma}[The square wave]\label{lem:square-wave}
Let $H\ge1$. There are $b_\nu$ $(1\le|\nu|\le H)$ with $b_\nu=0$ for even $\nu$ and $|b_\nu|\le\frac2{\pi|\nu|}$
such that, with
\[
T_H(y)=\sum_{1\le|\nu|\le H}b_\nu e(\nu y),\qquad E_H(y)=\frac2{H+1}\sum_{\substack{|\nu|\le H\\ \nu\text{ even}}}
\Bigl(1-\frac{|\nu|}{H+1}\Bigr)e(\nu y),
\]
we have $E_H\ge0$ and $|\sq(y)-T_H(y)|\le E_H(y)$ for all $y\notin\frac12\Z$.
\end{lemma}

\begin{proof}
For $y\notin\frac12\Z$ we have $\sq(y)=2\dpar{y+\frac12}-2\dpar{y}$. Indeed, put $f=\{y\}\ne0,\frac12$. If $f<\frac12$, then $\{y+\frac12\}=f+\frac12$ and $\dpar{y+\frac12}-\dpar{y}=f-(f-\frac12)=\frac12$; if $f>\frac12$, then $\{y+\frac12\}=f-\frac12$ and the difference is $(f-1)-(f-\frac12)=-\frac12$. These are $\frac12\sq(y)$ in the two cases.
Applying Lemma~\ref{lem:vaaler} at $y$ and at
$y+\frac12$, and using $e(\nu(y+\frac12))=(-1)^\nu e(\nu y)$, gives the claim with $b_\nu=2a_\nu((-1)^\nu-1)$ and error $2\Delta_H(y)+2\Delta_H(y+\frac12)=E_H(y)$,
which is $\ge0$.
\end{proof}

\begin{lemma}[Heads]\label{lem:heads}
Let $9\le K\le p/2$, let $B\subset[K,2K)$ be an interval, and let $H,N$ be integers with $1\le H<p$ and $2\le N<K$.
Then
\[
\Big|\sum_{\substack{k\in B\\ k\text{ odd}}}\frac{\Phi_k^{\le N}}{k}\Big|\le(1+\log N)\Bigl(\frac2H+10\,(1+\log H)\,
\frac{\mathcal{M}_K}{K}\Bigr),
\]
where $\mathcal{M}_K$ is defined in~\eqref{eq:MK}.
\end{lemma}

\begin{proof}
The idea is to replace $\sq$ by its truncated Fourier series (Lemma~\ref{lem:square-wave}); the resulting exponential sums over $k$ are controlled by $\mathcal{M}_K$ after Abel summation, and the error term of the approximation is a Fej\'er-type average of the same sums.

\emph{Step 1: expansion.} Put $y_{kn}=n\overline{k}/2p$, $\varphi_n(x)=h(n/x)/x$ and $w_{kn}=c_n(k)/k=\frac2{\pi n}\varphi_n(k)>0$ for odd
$k\in B$ and odd $n\le N$ (note $n<K\le k$). By \eqref{eq:Phik-sum} the sum to be estimated is
$\sum_{k,n}w_{kn}\sq(y_{kn})$. Let $Z(\nu)=\sum_{k,n}w_{kn}e(\nu y_{kn})$. By Lemma~\ref{lem:square-wave} and $w_{kn}>0$,
\begin{equation}\label{eq:heads-expand}
\Big|\sum_{k,n}w_{kn}\sq(y_{kn})-\sum_\nu b_\nu Z(\nu)\Big|\le\sum_{k,n}w_{kn}E_H(y_{kn})\le\frac2{H+1}\Bigl(Z(0)+
\sum_{\substack{0<|\nu|\le H\\ \nu\text{ even}}}|Z(\nu)|\Bigr).
\end{equation}

\emph{Step 2: the terms $\nu=0$ and $\nu\ne0$.} The number of odd $k\in B$ is at most $\frac K2+1$, so $\sum_{k\in B\text{ odd}}\frac1k\le\frac12+\frac1K\le1$. By the argument of~\eqref{eq:head-trivial},
$\sum_{n\text{ odd},\,n\le N}c_n(k)\le\frac2\pi(1+\frac12\log N)$, hence
\[
Z(0)\le\frac2\pi\Bigl(1+\frac12\log N\Bigr)\le1+\log N.
\]
For $\nu\ne0$ with $|\nu|\le H<p$ we have $p\nmid\nu n$, since $n<K\le p/2$. On $[K,2K]$ we have $0<\varphi_n\le1/K$ and
\[
|\varphi_n'(x)|=\Bigl|\frac n{x^3}h'\Bigl(\frac nx\Bigr)+\frac{h(n/x)}{x^2}\Bigr|\le\frac{1+\pi^2/4}{x^2},
\]
because $n/x\le1$ and $|h'|\le\pi^2/4$ on $[0,1]$: the power series of $1-h$ has positive coefficients, so $-h'$ is increasing on $[0,1]$, and $-h'(1)=\pi^2/4$. Partial summation over the odd $k\in B$, with the sums of $e(\nu n\overline{k}/2p)$ over sub-intervals of $B$ bounded by $\mathcal{M}_K$, gives
\begin{align*}
\Big|\sum_k\varphi_n(k)e\Bigl(\frac{\nu n\overline{k}}{2p}\Bigr)\Big|&\le\mathcal{M}_K\Bigl(\frac1K+\int_K^{2K}\frac{1+\pi^2/4}{x^2}\,dx\Bigr)\\
&=\mathcal{M}_K\Bigl(\frac1K+\frac{1+\pi^2/4}{2K}\Bigr)\le\frac{3\mathcal{M}_K}K.
\end{align*}
Summing over $n$ with the weights $\frac2{\pi n}$, as for $Z(0)$, we get $|Z(\nu)|\le3(1+\log N)\mathcal{M}_K/K$.

\emph{Step 3: conclusion.} We have $\sum_{\nu\text{ odd},\,|\nu|\le H}|b_\nu|\le\frac4\pi(1+\frac12\log H)$, and there are at most $H$ even $\nu$ with $0<|\nu|\le H$. By~\eqref{eq:heads-expand}, the quantity $|\sum_{k,n}w_{kn}\sq(y_{kn})|$ to be estimated is at most $|\sum_\nu b_\nu Z(\nu)|$ plus the right side of~\eqref{eq:heads-expand}, hence at most
\begin{align*}
&(1+\log N)\Bigl[\frac2{H+1}\Bigl(1+3H\frac{\mathcal{M}_K}K\Bigr)+\frac4\pi\Bigl(1+\frac12\log H\Bigr)\frac{3\mathcal{M}_K}K\Bigr]\\
&\qquad\le(1+\log N)\Bigl[\frac2{H+1}+\Bigl(6+\frac{12}\pi+\frac6\pi\log H\Bigr)\frac{\mathcal{M}_K}{K}\Bigr].
\end{align*}
Since $6+\frac{12}\pi<10$ and $\frac6\pi<10$, this implies the lemma.
\end{proof}

\begin{lemma}[Tails]\label{lem:tails}
Let $k$ be odd with $3\le k<p$ and let $2\le N<k$. For $N\le t<k$ put $G_k(t)=\sum_{n\text{ odd},\,N<n\le t}\sq(n\theta_k)$
and $D_k=\sup_{N\le t<k}|G_k(t)|/t$. Then \textup{(a)} $|\Phi_k^{>N}|\le\frac\pi2D_k\log\frac kN$, and
\textup{(b)} $|\Phi_k^{>N}|\le\frac1\pi\log k$.
\end{lemma}

\begin{proof}
(a) By \eqref{eq:Phik-sum}, $\Phi_k^{>N}=-\sum_{N<n<k}c(n)\sq(n\theta_k)$, where $c(t)=\frac1k\cot\frac{\pi t}{2k}$ is
decreasing on $[N,k]$ with $c(k)=0$. Partial summation, $G_k(N)=0$ and $c(k)=0$ give
$\Phi_k^{>N}=\int_N^kG_k(t)c'(t)\,dt$. Since $\sin u\ge\frac{2u}\pi$ on $[0,\frac\pi2]$,
$|c'(t)|=\frac\pi{2k^2}\csc^2\frac{\pi t}{2k}\le\frac\pi{2t^2}$, so
\[
|\Phi_k^{>N}|\le\int_N^kD_kt\cdot\frac\pi{2t^2}\,dt=\frac\pi2D_k\log\frac kN.
\]
(b) follows from $c_n(k)\le\frac2{\pi n}$ and $\sum_{n\text{ odd},\,3\le n<k}\frac1n\le\frac12\log k$ (here $n\ge3$ because $N\ge2$).
\end{proof}

\begin{lemma}[Generic $\theta$]\label{lem:generic}
Let $M\ge1$ and $N\ge2$ be integers and $0<\eta\le\frac12$, and let $\theta\in\R$ satisfy $\|2m\theta\|\ge\eta$ for
$1\le m\le M$. Let $\mathcal{S}$ be a set of consecutive odd integers in $(N,t]$, $t\ge N$, with
$n\theta\notin\frac12\Z$ for $n\in\mathcal{S}$. Then
\[
\frac1t\Big|\sum_{n\in\mathcal{S}}\sq(n\theta)\Big|\le\frac1M+\frac{2+\log M}{\eta N}.
\]
In particular, $D_k\le\frac1M+\frac{2+\log M}{\eta N}$ if $\|2m\theta_k\|\ge\eta$ for $1\le m\le M$.
\end{lemma}

\begin{proof}
Write $\mathcal{S}=\{n_0+2j:0\le j<|\mathcal{S}|\}$. For $1\le|\nu|\le M$,
\[
\Big|\sum_{n\in\mathcal{S}}e(\nu n\theta)\Big|=\Big|\sum_je(2\nu j\theta)\Big|\le\frac1{2\|2\nu\theta\|}\le\frac1{2\eta}.
\]
Lemma~\ref{lem:square-wave} with $H=M$ gives, using $\sum_{n\text{ odd}\le M}\frac1n\le1+\frac12\log M$ and that there are at most $M$ even $\nu$ with $0<|\nu|\le M$,
\[
\Big|\sum_{n\in\mathcal{S}}\sq(n\theta)\Big|\le\sum_{\nu\text{ odd}}\frac2{\pi|\nu|}\cdot\frac1{2\eta}+\frac2{M+1}
\Bigl(|\mathcal{S}|+\frac M{2\eta}\Bigr)\le\frac{2|\mathcal{S}|}{M+1}+\frac1\eta\Bigl(1+\frac2\pi+\frac1\pi\log M\Bigr).
\]
Now $|\mathcal{S}|\le\frac{t-N}2+1\le\frac t2$ since $N\ge2$, and $t\ge N$, so dividing by $t$ gives
\[
\frac1t\Big|\sum_{n\in\mathcal{S}}\sq(n\theta)\Big|\le\frac1{M+1}+\frac{1}{\eta N}\Bigl(1+\frac2\pi+\frac1\pi\log M\Bigr)\le\frac1M+\frac{2+\log M}{\eta N}.
\]
For the last statement take
$\mathcal{S}=\{n\text{ odd}:N<n\le t\}$ with $t<k$; by the remark after \eqref{eq:Phik-sum},
$n\theta_k\notin\frac12\Z$.
\end{proof}

\begin{lemma}[Witnesses]\label{lem:witnesses}
Let $k$ be odd with $3\le k<p$, let $1\le m<k$ and $0<\eta\le\frac12$. If $\|2m\theta_k\|<\eta$, then there is an
integer $r$ with $0<|r|<\eta k$ and $k\mid pr-m$.
\end{lemma}

\begin{proof}
Choose $b\in\Z$ with $|2mu_k/k-b|<\eta$ and put $r=2mu_k-bk$; then $|r|<\eta k$ and $pr\equiv2pu_k\,m\equiv m\pmod
k$. If $r=0$ then $k\mid2mu_k$; since $\gcd(2u_k,k)=1$ (because $2pu_k\equiv1\pmod k$), this gives $k\mid m$, which is impossible for $1\le m<k$.
\end{proof}

\section{Large divisors of the terms of an arithmetic progression}\label{sec:divisors}

This section proves Theorem~\ref{thm:divisors}. It is self-contained and does not involve $S(p)$. In Section~\ref{sec:proof} it is applied with $q=p$ and $T=\exp((\log p)^{1/20})$.

For $q\ge3$ put $L=\log q$ and $T=\exp(L^{\gamma_0})$. The theorem bounds the $1/d$-weighted mass of the divisors $d>TR$ of the $R$ terms of a progression $a+qt$ of moduli $q$ and size at most $q^2$. For a random set of $R$ integers of size at most $q^2$ this mass would be of order $1/T$, because a random integer $n$ has $\sum_{d\mid n,\,d>Y}1/d$ of average size $\sum_{d>Y}d^{-2}\approx1/Y$, and $Y=TR$. We prove the bound $T^{-\gamma_0/4}$ for every fixed $\gamma_0\in(0,1)$ and for every length $R\ge1$, with no assumption that $R$ is large compared with $q$; this is the difference from Shiu's theorem~\cite{b24}, which concerns progressions of length at least $q^\eta$ for a fixed $\eta>0$ and gives an upper bound of the correct order of magnitude for sums of multiplicative functions, not a power saving for the divisors beyond $TR$. We do not know whether Theorem~\ref{thm:divisors} follows from the results on divisors in short intervals and progressions in Hall and Tenenbaum~\cite{b37} or in Nair and Tenenbaum~\cite{b38}. The method is not new: writing a divisor as a smooth part times a rough part with few prime factors, and counting the rough parts by a binomial bound, is the device of the proofs of Shiu's theorem and of the theorem of Nair and Tenenbaum, and it goes back to Erd\H{o}s~\cite{b10}. Those theorems, and Henriot's version that is uniform in the discriminant~\cite{b15}, need the length to be at least a fixed power of the size of the terms. What is new here is the iteration in Lemma~\ref{lem:one-step}, which removes this condition at the price of a very large~$T$.

The idea is to write a large divisor $d$ of a term as $d=d_1d_2$, where $d_1$ is the part of $d$ composed of primes up to some $y$. If $d_1$ is at most the length $R$ of the progression, then the terms divisible by $d_1$ are regularly spaced, and the rough part $d_2$, composed of large primes, can be chosen in few ways. If $d_1>R$, then $d_1$ has a divisor $s$ slightly below $R$, and the terms divisible by $s$ form a much shorter progression (Lemma~\ref{lem:one-step}). After a bounded number of such steps the progressions are so short that Rankin's trick (Lemma~\ref{lem:rankin}) finishes the proof.

\subsection{Setup}

Throughout, $q\ge3$ is an integer (not necessarily prime), $L=\log q$ and $X=q^2$. For a nonzero integer $n$, ``$d\mid n$'' means that $d$ is a positive divisor of $|n|$. Theorem~\ref{thm:divisors} states that, for $0<\gamma_0<1$, $L\ge L_0(\gamma_0)$, $T=\exp(L^{\gamma_0})$, every integer $a$ coprime to $q$, every $R\ge1$ and every set $\mathcal I$ of at most $R$ consecutive integers with $0<|a+qt|\le q^2$ for $t\in \mathcal I$,
\begin{equation}\label{eq:div-restated}
\sum_{t\in \mathcal I}\ \sum_{\substack{d\mid a+qt\\ d>TR}}\frac1d\;\le\;T^{-\gamma_0/4}=\exp\Bigl(-\frac{\gamma_0}{4}L^{\gamma_0}\Bigr).
\end{equation}

An \emph{admissible progression} is a set $A=\{a+qt:t\in \mathcal I\}$, where $\mathcal I$ is a set of consecutive integers, $\gcd(a,q)=1$, and $0<|a+qt|\le X$ for $t\in \mathcal I$; its \emph{length} is $|\mathcal I|$. For $Y>0$ put
\[
\Delta(A;Y)=\sum_{t\in \mathcal I}\sum_{\substack{d\mid a+qt\\ d>Y}}\frac1d .
\]
Then the left side of \eqref{eq:div-restated} is $\Delta(A;TR)$, where $A$ is an admissible progression of length at most $R$. We use two simple facts.
\begin{enumerate}[label=(F\arabic*),leftmargin=*]
\item Every divisor of an element of $A$ is coprime to $q$. If $e\ge1$ is coprime to $q$, the $t\in \mathcal I$ with $e\mid a+qt$ form one residue class modulo $e$, so their number is at most $|\mathcal I|/e+1$.
\item If $Y\ge X$ then $\Delta(A;Y)=0$.
\end{enumerate}
For (F1): a prime dividing both $q$ and a divisor of $a+qt$ would divide $a$; and $e\mid a+qt$ is a congruence $qt\equiv-a\pmod e$ with $q$ invertible modulo $e$. For (F2): every divisor of $a+qt$ is at most $|a+qt|\le X$.

For $y\ge2$ let $M(y)=\prod_{\ell\le y}(1-1/\ell)^{-1}=\sum_{P^+(e)\le y}1/e$, the sum over all integers $e\ge1$ whose prime factors are at most $y$ ($\ell$ denotes a prime). By Rosser and Schoenfeld~\cite[(3.30)]{b23}, for $x>1$,
\begin{equation}\label{eq:RS330}
M(x)<e^\gamma(\log x)\Bigl(1+\frac1{\log^2x}\Bigr).
\end{equation}
For $y\ge e$ we have $\log y+1/\log y\le1+\log y$, so \eqref{eq:RS330} gives
\begin{equation}\label{eq:Mertens}
M(y)\le e^\gamma(1+\log y)\qquad(y\ge2).
\end{equation}
For $2\le y<e$ we have $M(y)=2<e^\gamma(1+\log2)=3.01\ldots$ directly.

\subsection{Three lemmas}

\begin{lemma}[Rankin]\label{lem:rankin}
For $Y>0$ and $0<\alpha<1$, $\Delta(A;Y)\le|\mathcal I|\,Y^{-\alpha}P(\alpha)$, where $P(\alpha)=\max_{1\le n\le X}\sigma_{\alpha-1}(n)$ and $\sigma_{\alpha-1}(n)=\sum_{d\mid n}d^{\alpha-1}$. If $0<\alpha\le\frac12$ and $L\ge3$, then
\[
\log P(\alpha)\le4(4L)^\alpha(\log\log(4L)+1).
\]
\end{lemma}

\begin{proof}
For $n\ne0$,
\[
\sum_{d\mid n,\,d>Y}\frac1d\le\sum_{d\mid n}\frac1d\Bigl(\frac dY\Bigr)^\alpha=Y^{-\alpha}\sigma_{\alpha-1}(|n|),
\]
which gives the first claim. For the second, $\sigma_{\alpha-1}(n)\le\prod_{\ell\mid n}(1-\ell^{\alpha-1})^{-1}$. Let $\ell_1<\dots<\ell_k$ be the prime factors of $n\le X$ and let $w_i$ be the $i$-th prime. Then $\ell_i\ge w_i$, and $w\mapsto(1-w^{\alpha-1})^{-1}$ is decreasing, so the product is at most $\prod_{i\le k}(1-w_i^{\alpha-1})^{-1}$. Moreover $\theta(w_k)=\log(w_1\cdots w_k)\le\log(\ell_1\cdots\ell_k)\le\log n\le2L$, where $\theta$ is Chebyshev's function. Since $\theta(\ell)\ge\ell/2$ for every prime $\ell\ge3$ (for $\ell\ge41$ by $\theta(x)>x(1-1/\log x)$ of Rosser and Schoenfeld~\cite{b23}, and for $\ell<41$ by inspection), we get $w_k\le4L$ (trivially if $w_k=2$). For $\alpha\le\frac12$ we have $\ell^{\alpha-1}\le2^{-1/2}$, and $-\log(1-x)\le x/(1-x)\le x/(1-2^{-1/2})\le4x$ for $0\le x\le2^{-1/2}$. Hence
\[
\log P(\alpha)\le4\sum_{\ell\le4L}\ell^{\alpha-1}\le4(4L)^\alpha\sum_{\ell\le4L}\frac1\ell\le4(4L)^\alpha(\log\log(4L)+1).
\]
For the last step, $\sum_{\ell\le x}1/\ell\le\sum_{\ell\le x}-\log(1-1/\ell)=\log M(x)$, and by \eqref{eq:RS330} this is at most $\gamma+\log\log x+1/\log^2x\le\log\log x+1$ for $x\ge12$; here $x=4L\ge12$.
\end{proof}

For $y\ge2$ put $N_y=\lfloor\log X/\log y\rfloor$ and, for $Z\ge1$,
\[
G_y(Z)=\sum_{0\le i\le\log Z/\log y}\binom{N_y}{i}.
\]

\begin{lemma}[One step]\label{lem:one-step}
Let $A$ be an admissible progression of length at most $R$, let $T'\ge1$ and $2\le y\le R$, and let $\mathcal{S}$ be the set of integers $s$ with $R/y<s\le R$ and $P^+(s)\le y$. For $s\in\mathcal{S}$ let $A_s=\{n/s:\ n\in A,\ s\mid n\}$. Then each nonempty $A_s$ is an admissible progression of length at most $R/s+1\le2R/s$, and
\[
\Delta(A;T'R)\le\frac{2M(y)}{T'}\sum_{j\ge0}2^{-j}G_y(2^{j+1}T'R)+\sum_{s\in\mathcal{S}}\frac1s\,\Delta\Bigl(A_s;\frac{T'R}{s}\Bigr).
\]
\end{lemma}

\begin{proof}
The terms with a small smooth part $d_1$ are counted by (F1) and the number of choices of the rough part; the terms with a large smooth part have a divisor $s\in\mathcal S$ and pass to $A_s$.

Write $A=\{a+qt:t\in \mathcal I\}$. \emph{The progressions $A_s$.} Let $s\in\mathcal{S}$ with $A_s\ne\emptyset$. By (F1), $\gcd(s,q)=1$ and the $t\in \mathcal I$ with $s\mid a+qt$ are $t=t_s+sv$, with $v$ running over a set of at most $|\mathcal I|/s+1$ consecutive integers. Then $(a+qt)/s=a'+qv$ with $a'=(a+qt_s)/s$, and $\gcd(a',q)=1$ because $a'+qv$ divides $a+qt$. The elements of $A_s$ are nonzero and at most $X/s\le X$ in absolute value. So $A_s$ is admissible of length at most $|\mathcal I|/s+1\le R/s+1\le2R/s$, as $s\le R$.

\emph{The decomposition.} Consider the pairs $(n,d)$ with $n\in A$, $d\mid n$ and $d>T'R$, and write $d=d_1d_2$, where $d_1$ is the largest $y$-smooth divisor of $d$ and $d_2$ is $y$-rough, that is, has no prime factor $\le y$. These pairs give $\Delta(A;T'R)=\sum_{(n,d)}1/d$. We split them according to whether $d_1\le R$ or $d_1>R$.

\emph{Case $d_1\le R$.} Fix $j\ge0$ and consider the pairs with $2^jT'R<d\le2^{j+1}T'R$. The map $(n,d)\mapsto(d_1,n,d_2)$ is injective. Given $d_1$ (which is coprime to $q$ by (F1)), there are at most $R/d_1+1\le2R/d_1$ elements $n$ with $d_1\mid n$, by (F1). Given $n$, the prime factors of $d_2$, counted with multiplicity, form a sub-multiset of size at most $\log d_2/\log y\le\log(2^{j+1}T'R)/\log y$ of the multiset of prime factors of $n$ that exceed $y$. That multiset has at most $N_y$ elements, because $y^{\Omega}\le|n|\le X$ if $\Omega$ is the number of prime factors of $n$ exceeding $y$; and a multiset with at most $N_y$ elements has at most $\binom{N_y}{i}$ sub-multisets of size $i$. So there are at most $G_y(2^{j+1}T'R)$ choices for $d_2$. Since $1/d<2^{-j}/(T'R)$ and $\sum_{P^+(d_1)\le y}1/d_1=M(y)$, these pairs contribute at most
\[
\sum_{j\ge0}\ \sum_{P^+(d_1)\le y}\frac{2R}{d_1}\cdot\frac{2^{-j}}{T'R}\,G_y(2^{j+1}T'R)=\frac{2M(y)}{T'}\sum_{j\ge0}2^{-j}G_y(2^{j+1}T'R).
\]

\emph{Case $d_1>R$.} List the prime factors of $d_1$, with multiplicity, in non-decreasing order, and let $s$ be the first partial product that exceeds $R/y$. It exists because $d_1>R\ge R/y$, and $s\le(R/y)\cdot y=R$ because the previous partial product is at most $R/y$ (the empty product is $1\le R/y$) and every prime factor is at most $y$. So $s\in\mathcal{S}$ and $s\mid d\mid n$. The map $(n,d)\mapsto(s,n/s,d/s)$ is injective, $d/s\mid n/s$, $d/s>T'R/s$, and $1/d=\frac1s\cdot\frac1{d/s}$. So these pairs contribute at most $\sum_{s\in\mathcal{S}}\frac1s\Delta(A_s;T'R/s)$.
\end{proof}

\begin{lemma}\label{lem:Gy}
Let $y\ge2$ and $Z\ge1$. Then $G_y(2^jZ)\le(1+N_y)^{1+j\log2/\log y}G_y(Z)$ for $j\ge0$. If $1+N_y\le y^{1/3}$, then
\[
\sum_{j\ge0}2^{-j}G_y(2^{j+1}Z)\le2.7\,(1+N_y)\,G_y(2Z).
\]
Moreover $\sum_{0\le i\le m}\binom Ni\le(eN/m)^m$ for all integers $N,m$ with $1\le m\le N$.
\end{lemma}

\begin{proof}
Write $N=N_y$. Since $\binom N{i+1}\le N\binom Ni$, we have $\sum_{i\le m+1}\binom Ni\le(1+N)\sum_{i\le m}\binom Ni$. Passing from $Z$ to $2^jZ$ raises the upper limit $\lfloor\log Z/\log y\rfloor$ by at most $\lceil j\log2/\log y\rceil\le1+j\log2/\log y$, which gives the first claim.

For the second claim, apply the first claim with $2Z$ in place of $Z$ and $(1+N)^{j\log2/\log y}\le(y^{1/3})^{j\log2/\log y}=2^{j/3}$. For $j\ge1$ this gives $G_y(2^{j+1}Z)\le(1+N)2^{j/3}G_y(2Z)$, while the term $j=0$ is $G_y(2Z)$. Hence the sum is at most
\[
G_y(2Z)\Bigl(1+(1+N)\sum_{j\ge1}2^{-2j/3}\Bigr)\le G_y(2Z)\bigl(1+1.7025(1+N)\bigr).
\]
If $N=0$ then $G_y\equiv1$ and the sum equals $2$. If $N\ge1$ then $1\le\frac12(1+N)$, and the sum is at most $2.21(1+N)G_y(2Z)$. In both cases the bound $2.7(1+N)G_y(2Z)$ holds.

The last claim is the standard bound
\[
\sum_{i\le m}\binom Ni\le\sum_{i\le m}\frac{N^i}{i!}\le\Bigl(\frac Nm\Bigr)^m\sum_{i\ge0}\frac{m^i}{i!}=\Bigl(\frac{eN}m\Bigr)^m ,
\]
using $(N/m)^i\le(N/m)^m$ for $i\le m\le N$.
\end{proof}

\subsection{Proof of Theorem~\ref{thm:divisors}}\label{subsec:proof-divisors}

The proof is an induction over a chain of progressions of decreasing length $R$. Put $\kappa=\log T=L^{\gamma_0}$. We measure the supremum of $\Delta(A;T_jR)$ over progressions of length at most $R$ by a quantity $b_j$; we bound it by the Rankin bound when $R$ is at most $T^\beta$, and by $6L\,b_{j+1}$ plus a tiny error when $R$ is larger, using Lemmas~\ref{lem:one-step} and~\ref{lem:Gy}. After at most $\lceil1/\gamma_0\rceil+1$ steps all lengths are below $T^\beta$.

\smallskip\noindent\emph{Overview.} For $j\ge0$ let $b_j$ be the supremum of $\Delta(A;2^{-j}TR)$ over admissible progressions $A$ of length at most $R$, for $R\le e^{x_j}$; we must show $b_0\le T^{-\gamma_0/4}$. Lemma~\ref{lem:one-step} splits a large divisor $d$ of a term into its $y$-smooth part $d_1$ and the rough rest $d_2$. If $d_1\le R$, the terms divisible by $d_1$ are regularly spaced and $d_2$ can be chosen in few ways (Lemma~\ref{lem:Gy}); this gives a tiny contribution $E_1=\exp(-\frac{2+\gamma_0}3\kappa+o(\kappa))$. If $d_1>R$, then $d_1$ has a divisor $s$ with $R/y<s\le R$, and the terms divisible by $s$ form an admissible progression of length at most $2R/s<2y$, so that $b_j\le E_b+E_1+6L\,b_{j+1}$, where $E_b$ is the base-case bound below. With $y=(2^{-j}TR)^{1/u}$ and $u\approx\kappa/(3\log L)$, the logarithm of the length falls from $x_j$ to $x_{j+1}=\log2+(\kappa+x_j)/u$, a contraction by a factor of about $u$, a power of $L$ up to logarithmic factors. Hence from $x_0=2L$ a bounded number $J\le j_0=\lceil1/\gamma_0\rceil+1$ of steps leads to $x_J\le\beta\kappa$, that is, to lengths at most $T^\beta$, where $\beta=\gamma_0/8$. For such lengths Rankin's trick (Lemma~\ref{lem:rankin}) with $\alpha=\gamma_0/2$ gives the base case $b_J\le E_b=\exp(-c_{\gamma_0}\kappa+o(\kappa))$ with $c_{\gamma_0}>\frac38\gamma_0$. Iterating the recursion at most $j_0$ times costs only a factor $(j_0+1)(6L)^{j_0}=e^{o(\kappa)}$, and the outcome is $b_0\le T^{-\gamma_0/4}$.

\medskip\noindent\emph{Step 0: what ``$L$ large'' means.} Write $\Lambda=\log L$, $j_0=\lceil1/\gamma_0\rceil+1$, $\alpha=\gamma_0/2$, $\beta=\gamma_0/8$ and
\[
c_{\gamma_0}=\tfrac38\gamma_0+\tfrac1{16}\gamma_0^2,\qquad
\eta_b=\alpha j_0\log2+4(4L)^{\alpha}(\log\log(4L)+1),\qquad
\eta_1=\frac\kappa\Lambda+2\Lambda+\log\bigl(32.4\cdot2^{j_0}\bigr).
\]
Let $L_0(\gamma_0)\ge3$ be so large that for all $L\ge L_0(\gamma_0)$ the following inequalities hold.
\begin{enumerate}[label=(L\arabic*),leftmargin=*]
\item $\kappa\ge93\Lambda$ and $\kappa\ge4(1+8/\gamma_0)\Lambda$.
\item $\beta\kappa\ge8\Lambda$.
\item $L^{\gamma_0}\ge(4\Lambda)^{j_0}$.
\item $j_0\log2\le\beta\kappa/2$.
\item $\eta_1\le\frac{11}{48}\kappa$.
\item $\log\bigl(2(j_0+1)\bigr)+j_0\log(6L)+\eta_b\le\frac{\gamma_0}8\kappa$.
\end{enumerate}
Each of these compares $\kappa=L^{\gamma_0}$ with a quantity that is $O_{\gamma_0}(L^{\gamma_0/2}\log\log L)$ or $O_{\gamma_0}(\log L)$ (for (L3), $L^{\gamma_0}$ with a power of $\log L$), so $L_0(\gamma_0)$ exists. Below, ``for $L\ge L_0$'' refers to these inequalities only; no other largeness is used. Since $\kappa=L^{\gamma_0}<L$ we have $\kappa\le L$, and $\Lambda\ge\log3>1$.

\medskip\noindent\emph{Step 1: the parameters.} Put
\[
u=\Bigl\lfloor\frac\kappa{3\Lambda}\Bigr\rfloor,\qquad R_b=T^\beta,\qquad T_j=2^{-j}T .
\]
Define $x_0=2L$, $x_{j+1}=\log2+(\kappa+x_j)/u$ and $J=\min\{j:\ x_j\le\beta\kappa\}$ (finiteness of $J$ is shown in (P3) below). For $j\le J$ and real $R\ge1$ let $B_j(R)$ be the supremum of $\Delta(A;T_jR)$ over all admissible progressions $A$ of length at most $R$ (the supremum of the empty set is $0$), and let $b_j=\sup\{B_j(R):1\le R\le e^{x_j}\}$. For $L\ge L_0$ we have the following.
\begin{enumerate}[label=(P\arabic*),leftmargin=*]
\item $u\ge\frac\kappa{3\Lambda}-1\ge\frac\kappa{4\Lambda}$ (as $\kappa\ge12\Lambda$), so by (L1) $u\ge1+8/\gamma_0$; also $u\ge93/3-1=30$. Since $\beta=\gamma_0/8$, $(u-1)\beta\ge1$. Also $\kappa/u\le\kappa/(\frac\kappa{3\Lambda}-1)=3\Lambda/(1-3\Lambda/\kappa)\le3.1\Lambda$, as $3\Lambda/\kappa\le1/31$.
\item The recursion $x_{j+1}=\log2+\kappa/u+x_j/u$ has the fixed point $x^*=(\log2+\kappa/u)\frac u{u-1}$ and $x_j-x^*=u^{-j}(x_0-x^*)$. Since $u\ge30$, $0<x^*\le(0.7+3.1\Lambda)\cdot\frac{30}{29}\le4\Lambda$ (as $\Lambda\ge1$), and $x_0-x^*\le x_0=2L$, so
\[
x_j\le4\Lambda+2Lu^{-j}.
\]
\item $J\le j_0$. Indeed, by (P1) and (L3), $u^{j_0}\ge(L^{\gamma_0}/(4\Lambda))^{j_0}=L^{\gamma_0j_0}/(4\Lambda)^{j_0}\ge L^{1+\gamma_0}/(4\Lambda)^{j_0}\ge L$, using $\gamma_0j_0\ge1+\gamma_0$. Hence $2Lu^{-j_0}\le2\le4\Lambda$, and by (P2) and (L2), $x_{j_0}\le8\Lambda\le\beta\kappa$. In particular $J$ is bounded in terms of $\gamma_0$ alone, and $2^J\le2^{j_0}$. Also $J\ge1$, since $x_0=2L>L>\beta\kappa$.
\item $J\log2\le\beta\kappa/2$, by (P3) and (L4).
\item $u\log2<\kappa/3$, because $u\le\kappa/(3\Lambda)$ and $\log2<\Lambda$.
\item $x_j\le2L$ for all $j$. By induction, if $x_j\le2L$ then $x_{j+1}\le\log2+(\kappa+2L)/u\le\log2+3L/30\le2L$, since $\kappa\le L$ and $u\ge30$.
\end{enumerate}
Note that $x_j>\beta\kappa$ for $j<J$ by the definition of $J$, so $e^{x_j}>R_b$ for $j<J$, while $e^{x_J}\le R_b$.

\medskip\noindent\emph{Step 2: the base case.} Let $j\le J$ and $1\le R\le R_b$. Since $\alpha=\gamma_0/2\le\frac12$ and $L\ge3$, Lemma~\ref{lem:rankin} applies with $Y=T_jR$ and $|\mathcal I|\le R$:
\[
B_j(R)\le R^{1-\alpha}T_j^{-\alpha}P(\alpha)\le R_b^{1-\alpha}\,2^{\alpha J}\,T^{-\alpha}P(\alpha),
\]
because $R^{1-\alpha}$ is increasing in $R$ and $T_j^{-\alpha}=2^{j\alpha}T^{-\alpha}\le2^{\alpha J}T^{-\alpha}$. Now $R_b^{1-\alpha}T^{-\alpha}=\exp((\beta(1-\alpha)-\alpha)\kappa)$ and
\[
\beta(1-\alpha)-\alpha=\frac{\gamma_0}8\Bigl(1-\frac{\gamma_0}2\Bigr)-\frac{\gamma_0}2=-\frac38\gamma_0-\frac1{16}\gamma_0^2=-c_{\gamma_0}.
\]
By (P3) $\alpha J\log2\le\alpha j_0\log2$, and by Lemma~\ref{lem:rankin} $\log P(\alpha)\le4(4L)^\alpha(\log\log4L+1)$. Hence
\begin{equation}\label{eq:div-Eb}
B_j(R)\le E_b:=\exp\bigl(-c_{\gamma_0}\kappa+\eta_b\bigr)\qquad(j\le J,\ 1\le R\le R_b),
\end{equation}
where $\eta_b=O_{\gamma_0}(L^{\gamma_0/2}\log\log L)=o(\kappa)$, so $E_b=\exp(-(\frac38\gamma_0+\frac1{16}\gamma_0^2)\kappa+o(\kappa))$. In particular $b_J\le E_b$, because $e^{x_J}\le R_b$.

\medskip\noindent\emph{Step 3: the recursive step.} Let $j<J$ and $R_b<R\le e^{x_j}$, and put $y=(T_jR)^{1/u}$, that is, $u\log y=\log(T_jR)$. Let $A$ be admissible of length at most $R$. We verify the hypotheses of Lemmas~\ref{lem:one-step} and~\ref{lem:Gy} and bound their output.
\begin{enumerate}[label=(\alph*),leftmargin=*]
\item \emph{$T_j\ge1$ and $y\le R$.} We have $\log T_j\ge\kappa-J\log2>0$ by (P4). Since $\log R>\beta\kappa$ and $(u-1)\beta\ge1$ by (P1), $\log T_j\le\kappa\le(u-1)\beta\kappa<(u-1)\log R$, so $\log y=(\log T_j+\log R)/u<\log R$.
\item \emph{Lower bound for $y$.} By (P4), $\log(T_jR)\ge\kappa-J\log2+\beta\kappa\ge(1+\frac\beta2)\kappa$. As $u\le\kappa/(3\Lambda)$,
\[
\log y\ge\Bigl(1+\frac\beta2\Bigr)\frac\kappa u\ge\Bigl(3+\frac{3\beta}2\Bigr)\Lambda\ge(3+\beta)\Lambda,
\]
so $y\ge L^{3+\beta}\ge27$ and, in particular, $y\ge2$.
\item \emph{$1+N_y\le y^{1/3}$.} Since $\log X=2L$, $N_y\le2L/\log y\le2L/(3\Lambda)$, so $1+N_y\le1+\frac23L\le L$ (as $L\ge3$). Also $y^{1/3}\ge L^{1+\beta/3}\ge L$. Thus $1+N_y\le L\le y^{1/3}$.
\item \emph{Upper limit in $G_y(2T_jR)$.} It is $\lfloor\log(2T_jR)/\log y\rfloor=\lfloor u+u\log2/\log(T_jR)\rfloor=u$, because $0<u\log2<\kappa\le\log(T_jR)$ by (P5).
\item \emph{The size of $G_y(2T_jR)$.} If $u\le N_y$, then $1\le u\le N_y$ and, using $u\log y=\log(T_jR)\ge\kappa$, $N_y/u\le2L/(u\log y)\le2L/\kappa=2L^{1-\gamma_0}$. Thus $eN_y/u\le2eL^{1-\gamma_0}$ and $eN_y/u\ge e>1$. Lemma~\ref{lem:Gy} with $m=u$, and $u\le\kappa/(3\Lambda)$, give
\[
G_y(2T_jR)=\sum_{i\le u}\binom{N_y}i\le\Bigl(\frac{eN_y}u\Bigr)^u\le\bigl(2eL^{1-\gamma_0}\bigr)^{\kappa/(3\Lambda)}=\exp\Bigl(\frac{1-\gamma_0}3\kappa+\frac{\log(2e)}{3\Lambda}\kappa\Bigr),
\]
where the second inequality uses that the base $2eL^{1-\gamma_0}$ is at least the base $eN_y/u\ge1$ and $u\le\kappa/(3\Lambda)$. If $u>N_y$, then $G_y(2T_jR)\le2^{N_y}\le2^u\le\exp(\kappa\log2/(3\Lambda))$ by (P5). Since $\log(2e)/3<1$ in both cases
\begin{equation}\label{eq:div-G}
G_y(2T_jR)\le\exp\Bigl(\frac{1-\gamma_0}3\kappa+\frac\kappa\Lambda\Bigr).
\end{equation}
\item \emph{$M(y)\le6L$.} By (P6) and $u\ge30$, $\log y\le(\log T_j+\log R)/u\le(\kappa+x_j)/u\le3L/30\le L$. Hence $M(y)\le e^\gamma(1+L)\le6L$ by \eqref{eq:Mertens}.
\end{enumerate}
By (a) and (b), Lemma~\ref{lem:one-step} applies to $A$ with $T'=T_j$, and by (c) and (d) Lemma~\ref{lem:Gy} applies with $Z=T_jR$. With \eqref{eq:div-G}, (c) and (e) we get
\[
\Delta(A;T_jR)\le\frac{2\cdot6L}{T_j}\cdot2.7L\cdot\exp\Bigl(\frac{1-\gamma_0}3\kappa+\frac\kappa\Lambda\Bigr)+\sum_{s\in\mathcal{S}}\frac1s\Delta\Bigl(A_s;\frac{T_jR}s\Bigr).
\]
For the first term, $T_j\ge2^{-J}T\ge2^{-j_0}T$ by (P3), and $12\cdot2.7=32.4$. So it is at most
\[
32.4\cdot2^{j_0}L^2\exp\Bigl(-\kappa+\frac{1-\gamma_0}3\kappa+\frac\kappa\Lambda\Bigr)=\exp\Bigl(-\frac{2+\gamma_0}3\kappa+\eta_1\Bigr)=:E_1,
\]
with $\eta_1=\kappa/\Lambda+2\Lambda+\log(32.4\cdot2^{j_0})=o(\kappa)$; thus $E_1=\exp(-\frac{2+\gamma_0}3\kappa+o(\kappa))$.

In the second term, $A_s$ has length at most $R':=2R/s$ by Lemma~\ref{lem:one-step}, and $T_jR/s=T_{j+1}R'$, because $T_{j+1}R'=\frac12T_j\cdot2R/s$. Moreover $2\le R'<2y$, because $R/y<s\le R$, and so, using $R\le e^{x_j}$,
\[
\log R'<\log2+\log y=\log2+\frac{\log T_j+\log R}u\le\log2+\frac{\kappa+x_j}u=x_{j+1}.
\]
Since $j+1\le J$, this gives $\Delta(A_s;T_jR/s)\le B_{j+1}(R')\le b_{j+1}$. As $\sum_{s\in\mathcal{S}}1/s\le M(y)\le6L$,
\begin{equation}\label{eq:div-Bj}
B_j(R)\le E_1+6L\,b_{j+1}\qquad(j<J,\ R_b<R\le e^{x_j}).
\end{equation}

\medskip\noindent\emph{Step 4: conclusion.} For $j<J$ and $1\le R\le e^{x_j}$ either $R\le R_b$ and $B_j(R)\le E_b$ by \eqref{eq:div-Eb}, or $R>R_b$ and \eqref{eq:div-Bj} applies. So $b_j\le E_b+E_1+6L\,b_{j+1}$ for $j<J$, and $b_J\le E_b$. Iterating, and using $6L\ge1$,
\[
b_0\le(E_b+E_1)\sum_{i<J}(6L)^i+(6L)^JE_b\le(J+1)(6L)^J(E_b+E_1)\le(j_0+1)(6L)^{j_0}(E_b+E_1).
\]
Next, $E_1\le E_b$: since $c_{\gamma_0}\le\frac38+\frac1{16}=\frac7{16}$ for $\gamma_0<1$, we have $\frac{2+\gamma_0}3-c_{\gamma_0}\ge\frac23-\frac7{16}=\frac{11}{48}$, so $E_1\le E_b$ follows from $\eta_1\le\frac{11}{48}\kappa+\eta_b$, which holds by (L5) and $\eta_b\ge0$. Therefore, by (L6) and $c_{\gamma_0}>\frac38\gamma_0$,
\[
\log b_0\le\log\bigl(2(j_0+1)\bigr)+j_0\log(6L)-c_{\gamma_0}\kappa+\eta_b\le-c_{\gamma_0}\kappa+\frac{\gamma_0}8\kappa\le-\frac{\gamma_0}4\kappa .
\]
That is, $b_0\le e^{-\gamma_0\kappa/4}=T^{-\gamma_0/4}$.

Finally let $a$, $R$ and $\mathcal I$ be as in the theorem, $A=\{a+qt:t\in \mathcal I\}$. If $\mathcal I=\emptyset$ the sum is zero. Otherwise $A$ is admissible of length at most $R$. If $R\ge X$, then $TR>X$ and $\Delta(A;TR)=0$ by (F2). Otherwise $1\le R<X$, so $\log R<2L=x_0$, and $\Delta(A;TR)\le B_0(R)\le b_0\le T^{-\gamma_0/4}$, since $T_0=T$. This is \eqref{eq:div-restated}. \qed

\begin{remark}\label{rem:optimality}
(a) \emph{Near-optimality.} Let $0<\varepsilon<\gamma_0$. The same proof with $\alpha=\gamma_0-\varepsilon/3$ and $\beta=\varepsilon/3$ gives the bound $T^{-\gamma_0+\varepsilon}$ for $L\ge L_0(\gamma_0,\varepsilon)$. The changes are as follows. The base case exponent is $\beta(1-\alpha)-\alpha=-\gamma_0+\frac23\varepsilon-\frac13\gamma_0\varepsilon+\frac19\varepsilon^2\le-\gamma_0+\varepsilon$ (since $\frac13\varepsilon+\frac13\gamma_0\varepsilon-\frac19\varepsilon^2\ge0$), and $\frac{2+\gamma_0}3>\gamma_0$ for $\gamma_0<1$, so the first term of the recursive step is smaller; if $\alpha>\frac12$ one uses $\log P(\alpha)\ll_\alpha L^\alpha/\log L$ in place of the bound of Lemma~\ref{lem:rankin} (the proof is the same, with $-\log(1-x)\le x/(1-2^{\alpha-1})$ for $0\le x\le 2^{\alpha-1}$, and the prime number theorem for $\sum_{\ell\le4L}\ell^{\alpha-1}$); the parameters $u$ and $x_j$ are unchanged, $J\le j_0$ because $\kappa=L^{\gamma_0}$ is unchanged, and (L1)--(L6) are adjusted to $\beta=\varepsilon/3$ (in (L1), $8/\gamma_0$ becomes $3/\varepsilon$).

This is essentially optimal, as the following example shows. Let $q$ be a prime, $R=1$, and let $z$ be maximal with $\theta(z)\le2L$, so that $z=(2+o(1))L$ and $P_z=\prod_{\ell\le z}\ell\le q^2$. For $L$ large $q>z$, so $a=P_z$ is coprime to $q$, and with $\mathcal I=\{0\}$ the progression is the single number $P_z$. Put $k=\lfloor\kappa/\log(z/2)\rfloor+1$. The products of $k$ distinct primes from $(z/2,z]$ are divisors of $P_z$ in $(e^{\kappa},z^k]$, since $(z/2)^k>e^\kappa=T$, and there are at least $\binom{\pi(z)-\pi(z/2)}k\ge\bigl((\pi(z)-\pi(z/2))/k\bigr)^k$ of them. Each has reciprocal at least $z^{-k}$, so
\[
\Delta(A;T)\ge\Bigl(\frac{\pi(z)-\pi(z/2)}{kz}\Bigr)^k .
\]
By the prime number theorem, $\pi(z)-\pi(z/2)=(1+o(1))L/\Lambda$, $z=(2+o(1))L$, $\log(z/2)=\Lambda+o(1)$, so $k=(\kappa/\Lambda)(1+o(1/\Lambda))$ (the rounding adds $O(1)=o(\kappa/\Lambda^2)$), and $kz/(\pi(z)-\pi(z/2))=(2+o(1))\kappa$. Since $\log\kappa=\gamma_0\Lambda$, this gives $\Delta(A;T)\ge\exp\bigl(-k(\log(2\kappa)+o(1))\bigr)=T^{-\gamma_0-(\log2+o(1))/\Lambda}$. Thus no bound $T^{-\gamma_0-\delta}$ with a fixed $\delta>0$ can hold for $T=\exp((\log q)^{\gamma_0})$, $R=1$ and prime $q$.


(b) The proof uses only (F1), (F2), Rankin's trick, Mertens' theorem in the form \eqref{eq:Mertens} and the inequality $\theta(\ell)\ge\ell/2$ for primes $\ell\ge3$; in particular it does not use Shiu's theorem~\cite{b24}. The large size of $T$ is essential for the argument: the base case needs $T^\alpha$ to beat $P(\alpha)=\exp(O(L^\alpha\log\log L))$, which is impossible if $T$ is a power of $\log q$. The threshold $L_0(\gamma_0)$ given by the proof is effective but has not been optimized.

(c) \emph{Size of the threshold.} The theorem is purely asymptotic. For $\gamma_0=\frac1{20}$, the value used in Section~\ref{sec:proof}, we have $j_0=21$, and condition (L3) reads $\Lambda\ge420\log(4\Lambda)$ with $\Lambda=\log\log q$; it holds only for $\log\log q>4073$. Below such sizes the inequality of Theorem~\ref{thm:divisors} can fail. For example, for $q=10^{12}+39$, $R=1$, $\mathcal I=\{0\}$ and $a=985496152350226952635200$, which is coprime to $q$ and at most $q^2$, the left-hand side is $5.155\ldots$, while $T^{-\gamma_0/4}=0.985\ldots$; the mechanism is that of the example in~(a). So the theorem, as proved, says nothing about values of $q$ within computational reach.
\end{remark}


\section{Proof of Theorem~\ref{thm:main}}\label{sec:proof}

Part~(a) of Theorem~\ref{thm:main} is Proposition~\ref{prop:upper}; for part~(b) we combine the reduction of Theorem~\ref{thm:decomp}(ii), which expresses $S(p)$ through $R_K(p)$ with $K=p^{3/4}(\log p)^4$, with the estimates of Sections~\ref{sec:kloosterman}--\ref{sec:divisors}. Writing $K_W$ for this cutoff, put
\[
\Sigma=R_{K_W}(p)=\sum_{\substack{3\le k<K_W\\ k\text{ odd}}}\frac{\Phi_k}{k}.
\]
For a parameter $N$ we split $\Phi_k=\Phi_k^{\le N}+\Phi_k^{>N}$ as in~\eqref{eq:Phik-sum}, calling $\Phi_k^{\le N}$ the \emph{head} and $\Phi_k^{>N}$ the \emph{tail} of $\Phi_k$, and $\Sigma$ is divided into four ranges, treated by four different tools, as in the table below, where $L=\log p$, $T=\exp(L^{1/20})$, $K_*=\exp(L^{9/10})$ and the entries are asymptotic as $p\to\infty$.

\medskip
\begin{center}
\small
\renewcommand{\arraystretch}{1.25}
\begin{tabular}{@{}p{3.0cm}p{6.0cm}p{4.2cm}@{}}
\toprule
Range & Tool & Contribution to $|\Sigma|$\\
\midrule
heads, $k<K_*$ & trivial bound~\eqref{eq:head-trivial} & $\dfrac{L^{19/20}}{2\pi}\,(1+o(1))$\\
heads, $K_*\le k<K_W$ & Lemma~\ref{lem:heads} and the Bourgain--Garaev bound (Corollary~\ref{cor:MK}) & $O\bigl(L^{7/10}(\log L)^4\bigr)$\\
tails, generic $k$ & Diophantine condition $\|2m\theta_k\|\ge1/T$; Lemmas~\ref{lem:generic} and~\ref{lem:tails}(a) & $O(1)$\\
tails, non-generic $k$ & Lemmas~\ref{lem:witnesses} and~\ref{lem:tails}(b); Theorem~\ref{thm:divisors} & $o(1)$\\
\bottomrule
\end{tabular}
\end{center}
\medskip

The first range dominates and determines the exponent $\frac{19}{20}$. The exponent is $\frac{19}{20}=\frac1{20}+\frac9{10}$: the head of each $\Phi_k$ costs $\log N\approx L^{1/20}$, and there are $\log K_*=L^{9/10}$ logarithmic scales below $K_*$. The scale $L^{9/10}$ is chosen so that the Bourgain--Garaev bound is needed only for $\log K\ge(\log p)^{9/10}$ (Section~\ref{sec:kloosterman}).

\subsection{Proof of Theorem~\ref{thm:main}(b)}

\begin{proof}
The proof splits $\Sigma$ at the cutoffs $N$ (head versus tail of each $\Phi_k$) and $K_*$ (trivial bound versus Bourgain--Garaev bound for the heads), and treats the tails by a dichotomy in the Diophantine behavior of $\theta_k$. Throughout, $p$ is a sufficiently large prime, $L=\log p$, and $o(1)$ denotes a quantity tending to $0$ as $p\to\infty$.

\smallskip
\emph{Step 1: parameters.} Put
\[
T=\exp\bigl(L^{1/20}\bigr),\quad\eta=\frac1T,\quad M=H=\lceil L^2\rceil,\quad N=\bigl\lceil TM(2+\log M)\bigr\rceil,
\]
\[
K_*=\exp\bigl(L^{9/10}\bigr),\qquad K_W=p^{3/4}L^4 .
\]
Then $T=\exp((\log p)^{\gamma_0})$ with $\gamma_0=\frac1{20}$, which is the form of $T$ in Theorem~\ref{thm:divisors} for $q=p$. Since $\log M=2\log L+O(L^{-2})$ and $N\le TM(2+\log M)+1$,
\begin{equation}\label{eq:logN}
\log N=L^{1/20}+O(\log L),\qquad\text{so}\qquad 1+\log N=(1+o(1))L^{1/20}.
\end{equation}
We claim that for $p$ large
\begin{equation}\label{eq:param-chain}
M<N<K_*<K_W<\frac p4 .
\end{equation}
Indeed, $M<N$ is clear. Next, $\log N=L^{1/20}+O(\log L)<L^{9/10}=\log K_*$ for large $L$. Further $\log K_W=\frac34L+4\log L>L^{9/10}$ for large $L$. Finally $K_W<p/4$ is equivalent to $4L^4<p^{1/4}$, true for large $p$. Also
\begin{equation}\label{eq:KW-range}
p^{3/4}(\log p)^4=K_W\le p^{9/10}
\end{equation}
for large $p$ (as $L^4\le p^{3/20}$), and by the choice of $N$
\begin{equation}\label{eq:etaN}
\eta N\ge M(2+\log M).
\end{equation}
We use the elementary estimates, for real $X\ge3$,
\begin{equation}\label{eq:sum-recip}
\sum_{\substack{k\text{ odd}\\ 3\le k<X}}\frac1k\le\frac12\log X,\qquad
\sum_{\substack{k\text{ odd}\\ 3\le k<X}}\frac{\log k}k\le1+\frac14(\log X)^2 .
\end{equation}
\emph{Proof of~\eqref{eq:sum-recip}.} The function $1/t$ is convex, so its value at the midpoint of $[k-1,k+1]$ is at most its mean over this interval: $\frac1k\le\frac12\int_{k-1}^{k+1}\frac{dt}t$. The intervals $[k-1,k+1]$ for odd $k\ge3$ are disjoint and, for $k<X$, lie in $[2,X+1]$; hence the first sum is at most $\frac12\log\frac{X+1}2\le\frac12\log X$. For the second, $t\mapsto\log t/t$ is decreasing on $[3,\infty)$ (its derivative is $(1-\log t)/t^2<0$ for $t>e$). Hence for odd $k\ge5$ we have $\frac{\log k}k\le\frac12\int_{k-2}^k\frac{\log t}t\,dt$. The intervals $[k-2,k]$ are disjoint and lie in $[3,X)$, so these terms sum to at most $\frac12\cdot\frac12\bigl((\log X)^2-(\log3)^2\bigr)\le\frac14(\log X)^2$. The term $k=3$ is $\frac{\log3}3<1$. \qed

\smallskip
\emph{Step 2: the reduction.} By~\eqref{eq:KW-range} the value $K=K_W$ is admissible in Theorem~\ref{thm:decomp}(ii), and~\eqref{eq:decomp-ii} gives
\begin{equation}\label{eq:S-Sigma}
S(p)=\frac p2\log p+c_0p+\frac{4p}{\pi}\,\Sigma+O\Bigl(\frac pL\Bigr),\qquad\Sigma=R_{K_W}(p)=\sum_{\substack{3\le k<K_W\\ k\text{ odd}}}\frac{\Phi_k}k .
\end{equation}
Write $\Sigma=\Sigma_{\mathrm{head}}+\Sigma_{\mathrm{tail}}$ with
\[
\Sigma_{\mathrm{head}}=\sum_{\substack{3\le k<K_W\\ k\text{ odd}}}\frac{\Phi_k^{\le N}}k,\qquad
\Sigma_{\mathrm{tail}}=\sum_{\substack{3\le k<K_W\\ k\text{ odd}}}\frac{\Phi_k^{>N}}k .
\]
Recall that $\Phi_k^{>N}=0$ for $k\le N+1$ (Eq.~\eqref{eq:Phik-sum}). Further split $\Sigma_{\mathrm{head}}=\Sigma_{\mathrm{head}}^{<K_*}+\Sigma_{\mathrm{head}}^{\ge K_*}$ according to $k<K_*$ or $k\ge K_*$.

\smallskip
\emph{Step 3: heads with $k<K_*$.} By~\eqref{eq:head-trivial}, $|\Phi_k^{\le N}|\le\frac2\pi(1+\frac12\log N)=\frac1\pi(2+\log N)$, and by the first estimate in~\eqref{eq:sum-recip} (note $K_*\ge3$),
\[
\bigl|\Sigma_{\mathrm{head}}^{<K_*}\bigr|\le\frac1\pi(2+\log N)\cdot\frac12\log K_*=\frac1{2\pi}(2+\log N)L^{9/10}.
\]
By~\eqref{eq:logN},
\begin{equation}\label{eq:head-low}
\bigl|\Sigma_{\mathrm{head}}^{<K_*}\bigr|\le\Bigl(\frac1{2\pi}+o(1)\Bigr)L^{19/20}.
\end{equation}

\smallskip
\emph{Step 4: heads with $K_*\le k<K_W$.} The idea is to apply Lemma~\ref{lem:heads} on dyadic blocks, where Corollary~\ref{cor:MK} supplies a saving of $(\log K)^{-3/2}$ in $\mathcal M_K/K$; the number of blocks is $O(L)$, and the saving at the bottom of the range, $(L^{9/10})^{-3/2}=L^{-27/20}$, is what remains.

Put $K_j=2^jK_*$ and $B_j=[K_j,2K_j)\cap[K_*,K_W)$ for $j\ge0$. The odd $k$ with $K_*\le k<K_W$ are partitioned by the $B_j$, and $B_j\ne\varnothing$ only if $K_j<K_W$, that is $j<\log_2(K_W/K_*)\le\log_2K_W=(\frac34L+4\log L)/\log2<1.1\,L$ for large $L$. So at most $1.1L+1\le2L$ blocks are nonempty. Fix a nonempty $B_j$. By~\eqref{eq:param-chain}, $9\le K_*\le K_j<K_W<p/4$ and $N<K_*\le K_j$, and $2\le N$, $1\le H<p$. Hence Lemma~\ref{lem:heads} applies to $K=K_j$, $B=B_j$ and gives
\[
\Bigl|\sum_{\substack{k\in B_j\\ k\text{ odd}}}\frac{\Phi_k^{\le N}}k\Bigr|\le(1+\log N)\Bigl(\frac2H+10(1+\log H)\frac{\mathcal M_{K_j}}{K_j}\Bigr).
\]
Since $K_j\le p/2$ and $\log K_j\ge\log K_*=L^{9/10}$, Corollary~\ref{cor:MK} gives
\begin{align*}
\frac{\mathcal M_{K_j}}{K_j}&\le16\sqrt2\,C_{BG}\,(\log L)^3\,L\,(\log K_j)^{-3/2}\\
&\le16\sqrt2\,C_{BG}\,(\log L)^3\,L\cdot L^{-27/20}=16\sqrt2\,C_{BG}\,(\log L)^3L^{-7/20}.
\end{align*}
Also $H=\lceil L^2\rceil\ge L^2$ and $1+\log H\le3\log L$ for large $L$. Summing over at most $2L$ blocks,
\begin{align*}
\bigl|\Sigma_{\mathrm{head}}^{\ge K_*}\bigr|
&\le(1+\log N)\cdot2L\Bigl(\frac2{L^2}+10\cdot3\log L\cdot16\sqrt2\,C_{BG}(\log L)^3L^{-7/20}\Bigr)\\
&=(1+\log N)\Bigl(\frac4L+960\sqrt2\,C_{BG}\,(\log L)^4L^{13/20}\Bigr).
\end{align*}
With $1+\log N\le2L^{1/20}$ from~\eqref{eq:logN} this gives
\begin{equation}\label{eq:head-high}
\bigl|\Sigma_{\mathrm{head}}^{\ge K_*}\bigr|\le1920\sqrt2\,C_{BG}\,L^{7/10}(\log L)^4+o(1)=O\bigl(L^{7/10}(\log L)^4\bigr),
\end{equation}
because $\frac1{20}+\frac{13}{20}=\frac7{10}$.

\smallskip
\emph{Step 5: generic tails.} The idea is that if $\theta_k$ is far from all rationals with denominator at most $M$ in the sense of $\|2m\theta_k\|$, then the partial sums of $\sq(n\theta_k)$ equidistribute, so that $D_k$ is small. Call an odd $k$ with $N+1<k<K_W$ \emph{generic} if $\|2m\theta_k\|\ge\eta$ for all $1\le m\le M$. Since $\eta=1/T\le\frac12$, $N\ge2$, and $k<K_W<p$, Lemma~\ref{lem:generic} and~\eqref{eq:etaN} give for generic $k$
\[
D_k\le\frac1M+\frac{2+\log M}{\eta N}\le\frac1M+\frac1M=\frac2M .
\]
Lemma~\ref{lem:tails}(a) applies (as $2\le N<k$) and gives
\[
|\Phi_k^{>N}|\le\frac\pi2D_k\log\frac kN\le\frac\pi M\log k .
\]
By the second estimate in~\eqref{eq:sum-recip} with $X=K_W$, $\log K_W=\frac34L+4\log L$ and $M\ge L^2$, the generic $k$ contribute to $\Sigma_{\mathrm{tail}}$ at most
\begin{equation}\label{eq:generic}
\frac\pi M\Bigl(1+\frac14\log^2K_W\Bigr)\le\frac\pi{L^2}\Bigl(1+\frac14\Bigl(\frac34L+4\log L\Bigr)^2\Bigr)=\frac{9\pi}{64}+o(1)=O(1).
\end{equation}

\smallskip
\emph{Step 6: non-generic tails.} The idea is that a non-generic $k$ is a large divisor of one of the integers $pr\mp m$ with $r$ small, and such divisors have small total reciprocal mass by Theorem~\ref{thm:divisors}.

Let $\mathcal N$ be the set of odd $k$ with $N+1<k<K_W$ that are not generic, that is, $\|2m\theta_k\|<\eta$ for some $1\le m\le M$. By Lemma~\ref{lem:tails}(b), for $k\in\mathcal N$
\begin{equation}\label{eq:nongeneric-trivial}
|\Phi_k^{>N}|\le\frac1\pi\log k\le\frac L\pi .
\end{equation}
Let $k\in\mathcal N$ and take such an $m$. Then $1\le m\le M<N<k<p$, so Lemma~\ref{lem:witnesses} (with $0<\eta\le\frac12$) gives an integer $r$ with $0<|r|<\eta k=k/T$ and $k\mid pr-m$. Call $(m,r)$ a \emph{witness} for $k$. Then
\[
k>T|r|,\qquad 1\le|r|<\frac kT<\frac{K_W}T .
\]
Let $R=2^i$, $i\ge0$ integer, with $R<K_W/T$. There are at most $\log_2K_W+1\le2L$ such $R$. For $1\le m\le M$ define
\[
A^+_{m,R}=\{pt-m:\ t\in\Z,\ R\le t<2R\},\qquad A^-_{m,R}=\{pt+m:\ t\in\Z,\ R\le t<2R\}.
\]
\emph{Admissibility in Theorem~\ref{thm:divisors} with $q=p$.} Each set is $\{a+pt:t\in\mathcal I\}$ with $a=\mp m$ and $\mathcal I=\{t:R\le t<2R\}$, a set of exactly $R$ consecutive integers. We have $\gcd(a,p)=1$ because $1\le m\le M<p$. The elements are positive and at most $2pR+M<2pK_W<p^2$: indeed, for $t\ge1$, $pt-m\ge p-M>0$, and $R<K_W/T\le K_W$ and $K_W<p/4$. So $0<|a+pt|\le p^2$ for all $t\in\mathcal I$. Finally $\log p=L\ge L_0(\gamma_0)$ for large $p$, and $T=\exp((\log p)^{\gamma_0})$ with $\gamma_0=\frac1{20}$ exactly.

Let $k\in\mathcal N$ have a witness $(m,r)$ and let $R=2^i$ with $R\le|r|<2R$; then $R<K_W/T$ since $R\le|r|<K_W/T$. If $r>0$ then $|pr-m|=pr-m\in A^+_{m,R}$; if $r<0$ then $|pr-m|=p|r|+m\in A^-_{m,R}$. In both cases $k$ is a positive divisor of an element $n$ of $A^\pm_{m,R}$ with $k>T|r|\ge TR$. For a finite set $\mathcal A$ of positive integers and $Y>0$ put
\[
\Delta(\mathcal A;Y)=\sum_{n\in\mathcal A}\ \sum_{\substack{d\mid n\\ d>Y}}\frac1d .
\]
Assigning to each $k\in\mathcal N$ one witness and the corresponding $(m,R,\pm)$, we get
\[
\sum_{k\in\mathcal N}\frac1k\le\sum_{m=1}^M\ \sum_{\substack{R=2^i\\ R<K_W/T}}\bigl(\Delta(A^+_{m,R};TR)+\Delta(A^-_{m,R};TR)\bigr).
\]
Theorem~\ref{thm:divisors} (with $\gamma_0=\frac1{20}$, so $T^{-\gamma_0/4}=T^{-1/80}$) bounds each of the $\le2\cdot M\cdot2L$ terms by $T^{-1/80}=\exp(-L^{1/20}/80)$. Hence
\[
\sum_{k\in\mathcal N}\frac1k\le4ML\,\exp\Bigl(-\frac{L^{1/20}}{80}\Bigr).
\]
With~\eqref{eq:nongeneric-trivial} and $M\le L^2+1\le2L^2$, the non-generic $k$ contribute to $\Sigma_{\mathrm{tail}}$ at most
\begin{equation}\label{eq:nongeneric}
\frac L\pi\cdot4ML\,e^{-L^{1/20}/80}\le3L^4e^{-L^{1/20}/80}=\exp\Bigl(\log3+4\log L-\frac{L^{1/20}}{80}\Bigr)=o(1).
\end{equation}

\smallskip
\emph{Step 7: conclusion.} Every odd $k$ with $3\le k<K_W$ is either $\le N+1$ (then $\Phi_k^{>N}=0$), generic, or in $\mathcal N$. Adding~\eqref{eq:head-low},~\eqref{eq:head-high},~\eqref{eq:generic} and~\eqref{eq:nongeneric},
\[
|\Sigma|\le\Bigl(\frac1{2\pi}+o(1)\Bigr)L^{19/20}+O\bigl(L^{7/10}(\log L)^4\bigr)+O(1)=O\bigl(L^{19/20}\bigr).
\]
Insert this in~\eqref{eq:S-Sigma}. Since $|c_0|$ is an absolute constant and $p/L\le pL^{19/20}$,
\[
S(p)=\frac12\,p\log p+O\bigl(p(\log p)^{19/20}\bigr),
\]
which is Theorem~\ref{thm:main}(b). In fact the proof gives $\frac{4}{\pi}|\Sigma|\le(\frac2{\pi^2}+o(1))L^{19/20}$.
\end{proof}

\subsection{Proof of Theorem~\ref{thm:main}(a) and of the consequences}

\begin{proof}[Proof of Theorem~\ref{thm:main}(a)]
This is Proposition~\ref{prop:upper}: $|S(p)|\le\frac2{\pi^2}\,p(\log p+2)^2$ for every odd prime $p$.
\end{proof}

\begin{proof}[Proof of Conjecture~3.3 and of $S(p)>Bp$]
Let $C_1,C_2,\varepsilon>0$ be given. By Theorem~\ref{thm:main}(b) there are $C$ and $p_0$ with $S(p)\ge\frac12p\log p-Cp(\log p)^{19/20}$ for $p\ge p_0$. Since $\frac12\log p-C(\log p)^{19/20}\to\infty$, it exceeds $C_1$, and indeed any fixed $B$, for all large $p$; thus $S(p)>C_1p$ and $S(p)>Bp$ for all sufficiently large primes. For the upper half, Theorem~\ref{thm:main}(a) gives $S(p)\le\frac2{\pi^2}p(\log p+2)^2$, and $(\log p+2)^2\le\frac{C_2\pi^2}2p^{\varepsilon}$ for all large $p$, so $S(p)<C_2p^{1+\varepsilon}$. Together, $C_1p<S(p)<C_2p^{1+\varepsilon}$ for all sufficiently large primes $p$, which is Conjecture~3.3.
\end{proof}

\begin{remark}\label{rem:sharper}
The exponent $\frac{19}{20}$ comes from the choice of parameters. Using the Bourgain--Garaev bound throughout the range covered by its proof, together with a variant of Theorem~\ref{thm:divisors} for smaller $T$, the same method gives a smaller power of $\log p$; we do not pursue this here. In the computed range, $S(p)/p-\frac12\log p$ stays between $-2.2$ and $2.4$ (Section~\ref{sec:computations}).
\end{remark}

\subsection{Correspondence with the Lean formalization}

The Lean proof (Appendix~\ref{app:lean}) follows the route of this section, with the parameters $T=\exp(L^{1/20})$, $M=H=\lceil L^2\rceil$, $K_*=\lceil\exp(L^{9/10})\rceil$, $K_W=p^{3/4}L^4$ and about $2L$ dyadic blocks; the choice of $N$ differs slightly from the one above. It establishes the upper bound of Theorem~\ref{thm:main}(a) and the statement $S(p)/p\to\infty$, from which Conjecture~3.3 follows; it does not establish the explicit exponent $\frac{19}{20}$. The Bourgain--Garaev bound enters the Lean proof as a hypothesis. In Lean, $K_W$ is rounded down to $\lfloor p^{3/4}(\log p)^4\rfloor$, which can lie just below the range of Theorem~\ref{thm:decomp}(ii) as stated; the two versions of $R_{K}(p)$ differ by at most one term, of size $O(p\log p/K_W)$, which is negligible. The following table lists the corresponding declarations.

\begin{center}
\small
\renewcommand{\arraystretch}{1.2}
\begin{tabular}{@{}p{5.3cm}p{8.4cm}@{}}
\toprule
Step of the proof & Lean declaration\\
\midrule
Theorem~\ref{thm:decomp}(i): identity $S=\Psi+\mathcal T$; asymptotics of $\Psi$ with error $\le\varepsilon p$ & \raggedright\texttt{TheoremB.S\_eq\_Psi\_Phi}; \texttt{TheoremB.Psi\_asymp}\tabularnewline
Theorem~\ref{thm:decomp}(ii), $K=\lfloor p^{3/4}(\log p)^4\rfloor$, error $O(p)$ (unspecified $c$ in place of $c_0$) & \raggedright\texttt{TheoremB.HBweak34\_of\_HK34}, \texttt{ThetaSum.hBweak34}\tabularnewline
Theorem~\ref{thm:decomp}(ii), $K=\lfloor p^{3/4}(\log p)^4\rfloor$, constant $c_0$, error $\le\varepsilon p$ & \raggedright\texttt{ThetaSum.theorem\_B\_c0}\tabularnewline
Lemma~\ref{lem:kloosterman} (Kloosterman's bound, fourth moment) & \raggedright\texttt{ThetaSum.hK34}, \texttt{kl\_four\_le}\tabularnewline
Heads (Step~4) & \raggedright\texttt{block\_heads}, \texttt{block\_kloost}\tabularnewline
Tails (Steps~5 and~6) & \raggedright\texttt{tails\_sum\_explicit}, \texttt{generic\_bound}, \texttt{nongeneric\_mass\_D}\tabularnewline
Theorem~\ref{thm:divisors} & \raggedright\texttt{ThetaSum.theorem\_D}\tabularnewline
Assembly (Steps~2--7) and Conjecture~3.3 & \raggedright\texttt{lower\_half\_gen}, \texttt{lower\_half\_HBweak34\_HBG}, \texttt{conjecture\_3\_3\_of\_BG}\tabularnewline
\bottomrule
\end{tabular}
\end{center}

\section{Computations}\label{sec:computations}

This section proves the counting formula that underlies the computation of $S(p)$, describes the algorithm and its checks, and proves Proposition~\ref{prop:numerics}. It then compares the computed values with the asymptotic formula of Theorem~\ref{thm:main}. The data and programs are available at \url{https://github.com/samsonmaximus/theta-sum}.

\subsection{A counting formula}
For integers $x\ge0$ and $y\ge1$ put
\[
W(x,y)=\sum_{c=1}^{\lfloor(y-1)/2\rfloor}\Bigl\lfloor\frac{cx}{y}\Bigr\rfloor .
\]

\begin{lemma}[Counting formula]\label{lem:counting}
For odd $a$ with $1\le a<p$,
\[
U(a)=(p-1)a-4W(2p,a),\qquad\text{hence}\qquad S(p)=\frac{(p-1)^3}4-4\sum_{\substack{3\le a\le p-2\\ a\text{ odd}}}W(2p,a).
\]
Moreover, for integers $q,r\ge0$ and $y\ge1$, with $A=\lfloor(y-1)/2\rfloor$,
\begin{equation}\label{eq:Lambda-shift}
W(qy+r,y)=q\binom{A+1}2+W(r,y),
\end{equation}
and for coprime integers $r,y\ge1$
\begin{equation}\label{eq:Lambda-recip}
W(r,y)+W(y,r)=\Bigl\lfloor\frac{y-1}2\Bigr\rfloor\Bigl\lfloor\frac{r-1}2\Bigr\rfloor .
\end{equation}
\end{lemma}

\begin{proof}
The idea is to count, for each $c$, the $b$ with $\lfloor ab/p\rfloor=c$; this turns $U(a)$ into an alternating sum of $\lfloor cp/a\rfloor$, of which only the even $c$ need be kept. The last identity is a count of lattice points on the two sides of a line.

\emph{Step 1: the formula for $U(a)$.} Let $a$ be odd and $U(a)=\sum_{b=1}^{p-1}(-1)^{\lfloor ab/p\rfloor}$ as in Lemma~\ref{lem:elementary}. For $0\le c\le a$ let $M_c$ be the number of $b\in[1,p-1]$ with $ab<cp$. Then $M_0=0$ and $M_a=p-1$. For $1\le c\le a-1$ we have $M_c=\lfloor cp/a\rfloor$: the condition is $b<cp/a$, which implies $b<p$, and $a\nmid cp$ because $a$ is coprime to $p$ and $a\nmid c$. For $0\le c\le a-1$ the integer $\lfloor ab/p\rfloor$ equals $c$ for exactly $M_{c+1}-M_c$ values of $b$ (note that $ab\ne cp$ for $1\le b\le p-1$, $c\ge1$). Summation by parts gives, since $a$ is odd,
\[
U(a)=\sum_{c=0}^{a-1}(-1)^c(M_{c+1}-M_c)=-(-1)^aM_a-2\sum_{c=1}^{a-1}(-1)^cM_c=(p-1)-2\sum_{c=1}^{a-1}(-1)^c\Bigl\lfloor\frac{cp}a\Bigr\rfloor .
\]
Pairing $c$ with $a-c$ and using $\lfloor cp/a\rfloor+\lfloor(a-c)p/a\rfloor=p-1$ (as $cp/a\notin\Z$) gives $\sum_{c=1}^{a-1}\lfloor cp/a\rfloor=(p-1)(a-1)/2$. The terms with $c=2c'$ even sum to $\sum_{c'=1}^{(a-1)/2}\lfloor 2c'p/a\rfloor=W(2p,a)$, because $\lfloor(a-1)/2\rfloor=(a-1)/2$. Therefore
\[
\sum_{c=1}^{a-1}(-1)^c\Bigl\lfloor\frac{cp}a\Bigr\rfloor=2W(2p,a)-\frac{(p-1)(a-1)}2,
\]
and $U(a)=(p-1)+(p-1)(a-1)-4W(2p,a)=(p-1)a-4W(2p,a)$.

\emph{Step 2: the formula for $S(p)$.} By Lemma~\ref{lem:elementary}(b), $S(p)=\sum U(a)$ over odd $a$ with $1\le a\le p-2$. There are $(p-1)/2$ such $a$, and $\sum a=((p-1)/2)^2$, so $\sum(p-1)a=(p-1)^3/4$. Finally $W(2p,1)=0$.

\emph{Step 3: the two relations.} Identity \eqref{eq:Lambda-shift} follows from $\lfloor c(qy+r)/y\rfloor=cq+\lfloor cr/y\rfloor$ and $\sum_{c=1}^Ac=\binom{A+1}2$. For \eqref{eq:Lambda-recip} put $B=\lfloor(r-1)/2\rfloor$. For $1\le c\le A$ we have $cr/y<r/2$, and $y\nmid cr$ (otherwise $y\mid c$, as $\gcd(r,y)=1$). So $\lfloor cr/y\rfloor$ is the number of integers $d\ge1$ with $dy<cr$, and each such $d$ satisfies $d<r/2$, that is, $d\le B$. Thus $W(r,y)$ is the number of lattice points $(c,d)$ with $1\le c\le A$, $1\le d\le B$ and $dy<cr$. In the same way $W(y,r)=\sum_{d=1}^B\lfloor dy/r\rfloor$ is the number of those points with $cr<dy$. No point of $[1,A]\times[1,B]$ has $dy=cr$. Hence $W(r,y)+W(y,r)=AB$.
\end{proof}

\subsection{Algorithm and checks}
For odd $a<p$ the integers $2p$ and $a$ are coprime. For coprime $x\ge1$, $y\ge2$ write $x=qy+r$ with $0<r<y$ (here $r\ne0$ since $\gcd(x,y)=1$) and $A=\lfloor(y-1)/2\rfloor$, $B=\lfloor(r-1)/2\rfloor$. Then \eqref{eq:Lambda-shift} and \eqref{eq:Lambda-recip} give
\[
W(x,y)=q\binom{A+1}2+AB-W(y,r),
\]
and $W(x,1)=0$. The pair $(x,y)$ is thus replaced by $(y,r)$ exactly as in the Euclidean algorithm. Hence $W(2p,a)$ is evaluated with $O(\log a)$ arithmetic operations, and by Lemma~\ref{lem:counting}, $S(p)$ is evaluated with $O(p\log p)$ operations, in exact integer arithmetic.

The program (\texttt{code/thetasum.c}) is written in C and uses exact 128-bit integer arithmetic; there are 283145 odd primes $p\le4\cdot10^6$ and 129703 primes $4\cdot10^6<p\le6\cdot10^6$. The values $S(p)$ for all 412848 odd primes $p\le6\cdot10^6$ were compared with the following:
\begin{enumerate}[label=\textup{(\roman*)},leftmargin=*,nosep]
\item a brute-force evaluation of the definition of $S(p)$ for all 668 odd primes $p<5000$ (option \texttt{brute} of \texttt{code/thetasum.c});
\item a second implementation of the algorithm, for the 129703 primes $4\cdot10^6<p\le6\cdot10^6$ (\texttt{paper/checks/sp.c});
\item a program that shares no code or method with the first (\texttt{code/floorsum\_check.c}): writing $(-1)^x=1-2x+4\lfloor x/2\rfloor$, it evaluates each inner sum over $j$ in the definition of $S(p)$ as a floor sum of a linear function by a Euclid-type algorithm, with $O(p\log p)$ operations in all. It was run on all 17983 odd primes $p<2\cdot10^5$, and on 640 further values: primes chosen at random from all ranges up to $6\cdot10^6$, together with the 20 extreme values of $S(p)/p-\frac12\log p$ for $10^5<p\le5\cdot10^6$.
\end{enumerate}
All values agree. The outputs of~(iii) are in the files \texttt{data/crosscheck\_floorsum*.txt} of the repository.

\subsection{Proof of Proposition~\ref{prop:numerics}}
The proposition is a statement about the finitely many values $S(p)$, $p\le6\cdot10^6$, computed as above. From the list of values (file \texttt{S\_p\_upto\_6e6.txt.gz} in the repository), $S(p)>0$ for every odd prime $p\le6\cdot10^6$, and for $m=1,\dots,5$ the largest prime $p\le6\cdot10^6$ with $S(p)\le mp$ is
\[
m=1,\dots,5:\qquad p=5,\ 233,\ 3119,\ 58013,\ 869399 .
\]
Hence $S(p)>mp$ for all primes $p$ with $p_m<p\le6\cdot10^6$, where $p_1,\dots,p_5$ are these five primes; this gives the five inequalities of the proposition. None of the thresholds can be lowered, because $p_m$ itself violates the inequality:
\[
\begin{aligned}
S(5)&=4\le5, & S(233)&=440\le466, & S(3119)&=9194\le9357,\\
S(58013)&=218428<232052, & S(869399)&=4287874<4346995,
\end{aligned}
\]
where $466=2\cdot233$, $9357=3\cdot3119$, $232052=4\cdot58013$ and $4346995=5\cdot869399$.
These values are the ones given by the computation; the values $S(233)=440$ and $S(3119)=9194$ are also verified in Lean (Appendix~\ref{app:lean}). \qed

The minimum of $S(p)/p$ over $3119<p\le6\cdot10^6$ is $3.0254$, attained at $p=4139$; over $58013<p\le6\cdot10^6$ it is $4.0055$ (at $p=90089$), and over $869399<p\le6\cdot10^6$ it is $5.1332$ (at $p=1021019$). For $5\cdot10^6<p\le6\cdot10^6$ we have $S(p)/p>5.82$.

\subsection{Comparison with the asymptotic formula}
Put $X_p=S(p)/p-\frac12\log p$. By Theorem~\ref{thm:decomp}(i), which holds for every odd prime, $X_p=c_0+\Tpart(p)/p+O(1/p)$, with $c_0=0.23993\ldots$ as in \eqref{eq:c0}. (Part~(ii) of that theorem says nothing in the computed range, since it needs $p>10^{56}$.) Over $10^3<p\le6\cdot10^6$ (412681 primes) the mean of $X_p$ is $0.2396$, close to $c_0$, so that numerically $\Tpart(p)/p$ has mean close to $0$; we do not prove this. The standard deviation of $X_p$ is $0.626$; over all odd primes $p\le6\cdot10^6$, $X_p$ lies between $-2.160$ (at $p=750749$) and $2.386$ (at $p=4158001$). Table~\ref{tab:statistics} shows that the distribution is stable over the range.

\begin{table}[ht]
\small
\centering
\begin{tabular}{lrrrrr}
\toprule
range of $p$ & primes & mean & st.\ dev. & minimum (at $p$) & maximum (at $p$)\\
\midrule
$(10^3,\,2\cdot 10^5]$ & 17816 & 0.2387 & 0.624 & $-1.719$ (58013) & 2.219 (150151)\\
$(2\cdot 10^5,\,10^6]$ & 60514 & 0.2400 & 0.627 & $-2.160$ (750749) & 2.279 (910981)\\
$(10^6,\,2\cdot 10^6]$ & 70435 & 0.2397 & 0.625 & $-1.892$ (1767149) & 2.363 (1735651)\\
$(2\cdot 10^6,\,3\cdot 10^6]$ & 67883 & 0.2407 & 0.626 & $-2.041$ (2972969) & 2.320 (2238391)\\
$(3\cdot 10^6,\,4\cdot 10^6]$ & 66330 & 0.2391 & 0.625 & $-1.957$ (3950099) & 2.360 (3387931)\\
$(4\cdot 10^6,\,5\cdot 10^6]$ & 65367 & 0.2396 & 0.626 & $-1.840$ (4046543) & 2.386 (4158001)\\
$(5\cdot 10^6,\,6\cdot 10^6]$ & 64336 & 0.2390 & 0.628 & $-1.939$ (5525519) & 2.349 (5694151)\\
\bottomrule
\end{tabular}
\caption{Statistics of $S(p)/p-\frac12\log p$ over primes $p$ in each range.}\label{tab:statistics}
\end{table}

\subsection{The open range}
Conjectures~1.1, 3.1 and~3.2 hold for $p\le6\cdot10^6$ (Proposition~\ref{prop:numerics}) and for all $p$ larger than a threshold that has not been computed. The threshold would be effective if the implied constant in the Bourgain--Garaev bound were made explicit, but it is enormous: the proof of Theorem~\ref{thm:divisors} alone needs $\log\log p>4073$ (Remark~\ref{rem:optimality}(c)). Closing the gap between $6\cdot10^6$ and this threshold would require explicit constants in the Bourgain--Garaev bound (Lemma~\ref{lem:bg}), and an explicit form of the other estimates of Sections~\ref{sec:decomposition}--\ref{sec:proof}.

\section*{Use of AI}

This work is the result of human--AI interaction with Claude Opus~5.5 (Anthropic). I chose the problem and directed the project; the AI system found the approach, developed the proofs, and wrote the code, the Lean formalization and the text. Separate AI sessions refereed each proof, checked each citation against its source and recomputed every computational claim with independent code. I take full responsibility for the contents.

\begin{thebibliography}{99}

\bibitem{b1} B. C. Berndt, R. N. Bhat, J. L. Meyer, L. Xie and A. Zaharescu, \emph{An arithmetic sum associated with the classical theta function}, Results Math. \textbf{81} (2026), no.~1, Paper No.~25, doi:10.1007/s00025-025-02584-2; arXiv:2501.03234.

\bibitem{b2} B. C. Berndt, \emph{Analytic Eisenstein series, theta-functions, and series relations in the spirit of Ramanujan}, J. Reine Angew. Math. \textbf{303/304} (1978), 332--365.

\bibitem{b3} B. C. Berndt and L. A. Goldberg, \emph{Analytic properties of arithmetic sums arising in the theory of the classical theta-functions}, SIAM J. Math. Anal. \textbf{15} (1984), 143--150.

\bibitem{b4} S. Bettin, \emph{On the distribution of a cotangent sum}, Int. Math. Res. Not. IMRN \textbf{2015}, no.~21, 11419--11432.

\bibitem{b5} S. Bettin and J. B. Conrey, \emph{Period functions and cotangent sums}, Algebra Number Theory \textbf{7} (2013), 215--242.

\bibitem{b6} S. Bettin and S. Drappeau, \emph{Partial sums of the cotangent function}, J. Th\'eor. Nombres Bordeaux \textbf{32} (2020), 217--230.

\bibitem{b7} S. Bettin and S. Drappeau, \emph{Limit laws for rational continued fractions and value distribution of quantum modular forms}, Proc. Lond. Math. Soc. (3) \textbf{125} (2022), no.~6, 1377--1425.

\bibitem{b8} J. Bourgain and M. Z. Garaev, \emph{Sumsets of reciprocals in prime fields and multilinear Kloosterman sums}, Izv. Math. \textbf{78} (2014), no.~4, 656--707; Russian original: Izv. Ross. Akad. Nauk Ser. Mat. \textbf{78} (2014), no.~4, 19--72; arXiv:1211.4184. Section and page numbers refer to arXiv v1 (18 November 2012), the only arXiv version, unless stated otherwise.

\bibitem{b9} J. Bourgain and M. Z. Garaev, \emph{Kloosterman sums in residue rings}, Acta Arith. \textbf{164} (2014), 43--64.

\bibitem{b10} P. Erd\H{o}s, \emph{On the sum $\sum_{k=1}^{x}d(f(k))$}, J. London Math. Soc. \textbf{27} (1952), 7--15.

\bibitem{b11} Google DeepMind and contributors, \emph{formal-conjectures}, a repository of formalized mathematical conjectures in Lean, \url{https://github.com/google-deepmind/formal-conjectures}, file \texttt{FormalConjectures/Arxiv/2501.03234/ArithmeticSumS.lean} (accessed November 2026).

\bibitem{b12} M. Z. Garaev and I. E. Shparlinski, \emph{On the distribution of modular inverses from short intervals}, Mathematika \textbf{69} (2023), no.~4, 1183--1194.

\bibitem{b13} S. W. Graham and G. Kolesnik, \emph{Van der Corput's Method of Exponential Sums}, London Math. Soc. Lecture Note Ser. \textbf{126}, Cambridge University Press, 1991.

\bibitem{b14} G. H. Hardy, \emph{On certain series of discontinuous functions connected with the modular functions}, Quart. J. Pure Appl. Math. \textbf{36} (1905), 93--123.

\bibitem{b15} K. Henriot, \emph{Nair--Tenenbaum bounds uniform with respect to the discriminant}, Math. Proc. Cambridge Philos. Soc. \textbf{152} (2012), no.~3, 405--424; Erratum, ibid. \textbf{157} (2014), no.~2, 375--377.

\bibitem{b17} M. A. Korolev, \emph{Karatsuba's method for estimating Kloosterman sums}, Sb. Math. \textbf{207} (2016), no.~8, 1142--1158.

\bibitem{b18} A. L\"ageler, \emph{Continued fractions and Hardy sums}, Abh. Math. Semin. Univ. Hambg. \textbf{94} (2024), no.~2, 107--128.

\bibitem{b23} J. B. Rosser and L. Schoenfeld, \emph{Approximate formulas for some functions of prime numbers}, Illinois J. Math. \textbf{6} (1962), 64--94.

\bibitem{b24} P. Shiu, \emph{A Brun--Titchmarsh theorem for multiplicative functions}, J. Reine Angew. Math. \textbf{313} (1980), 161--170.

\bibitem{b25} I. E. Shparlinski, \emph{Modular hyperbolas}, Jpn. J. Math. \textbf{7} (2012), 235--294.

\bibitem{b26} R. Sitaramachandrarao, \emph{Dedekind and Hardy sums}, Acta Arith. \textbf{48} (1987), 325--340.

\bibitem{b27} J. D. Vaaler, \emph{Some extremal functions in Fourier analysis}, Bull. Amer. Math. Soc. (N.S.) \textbf{12} (1985), 183--216.

\bibitem{b34} The mathlib Community, \emph{The Lean mathematical library}, in: Proceedings of the 9th ACM SIGPLAN International Conference on Certified Programs and Proofs (CPP 2020), ACM, 2020, 367--381.

\bibitem{b35} L. de Moura and S. Ullrich, \emph{The Lean 4 theorem prover and programming language}, in: Automated Deduction -- CADE 28, Lecture Notes in Comput. Sci. \textbf{12699}, Springer, 2021, 625--635.

\bibitem{b36} H. D. Kloosterman, \emph{On the representation of numbers in the form $ax^2+by^2+cz^2+dt^2$}, Acta Math. \textbf{49} (1926), 407--464.

\bibitem{b37} R. R. Hall and G. Tenenbaum, \emph{Divisors}, Cambridge Tracts in Mathematics \textbf{90}, Cambridge University Press, 1988.

\bibitem{b38} M. Nair and G. Tenenbaum, \emph{Short sums of certain arithmetic functions}, Acta Math. \textbf{180} (1998), 119--144.

\end{thebibliography}

\appendix
% Environment leancode: typewriter block for Lean listings (spaces and line breaks kept; $...$ math works; inside math \sb and \sp print _ and ^).
\makeatletter
\newenvironment{leancode}{%
  \trivlist\item\relax
  \leftskip\@totalleftmargin\rightskip\z@skip
  \parindent\z@\parfillskip\@flushglue\parskip\z@skip
  \@@par
  \@tempswafalse
  \def\par{%
    \if@tempswa
      \leavevmode\null\@@par\penalty\interlinepenalty
    \else
      \@tempswatrue
      \ifhmode\@@par\penalty\interlinepenalty\fi
    \fi}%
  \ttfamily\small\frenchspacing\@noligs
  \catcode`\_=12 \catcode`\^=12
  \obeylines\@vobeyspaces}{\endtrivlist}
\makeatother

\section{The Lean formalization}\label{app:lean}

The formal proof follows the route of Sections~\ref{sec:notation}--\ref{sec:proof}. It is written in the Lean~4 proof assistant~\cite{b35} with the Mathlib library~\cite{b34} and contains no unproved steps; its main theorem takes one hypothesis, a form of the Bourgain--Garaev bound, which is not formalized. The project (Lean and Mathlib version \lean{v4.35.0}) has about 17000 lines and is in the repository \url{https://github.com/samsonmaximus/theta-sum}. It contains no \texttt{sorry} and no \lean{native_decide}. The build succeeds, and for each of the 56 distinct declarations printed in \lean{ThetaSum/MainTheorems.lean}, \texttt{\#print axioms} reports only \texttt{propext}, \lean{Classical.choice} and \lean{Quot.sound}.

Declaration names below lie in the namespace \lean{ThetaSum}, except \lean{theorem_A} and the names \lean{conjecture_4_*}, which lie in the namespace \lean{Arxiv.2501.03234} of~\cite{b11} (in Lean its last component is written in guillemets); names with the prefix \lean{TheoremB} or \lean{Recip} are written with their namespace.

\subsection*{Definitions}
The definitions $S'$ and $S$ are copied verbatim from the formal-conjectures file \lean{FormalConjectures/Arxiv/2501.03234/ArithmeticSumS.lean}~\cite{b11}:
\begin{leancode}
def S' (h k : $\N$) : $\Z$ :=
  $\sum$ j $\in$ Finset.Ico 1 k, (-1 : $\Z$) ^ (j + 1 + $\lfloor$(h * j : $\mathbb{Q}$) / k$\rfloor\sb{+}$)
def S (k : $\N$) : $\Z$ := $\sum$ h $\in$ Finset.Ico 1 k, S' h k
\end{leancode}

\subsection*{The hypotheses}
The hypothesis \texttt{HBG} is the following statement, in which $n^*$ is the inverse of $n$ modulo $p$.

\smallskip
\noindent\textbf{(HBG)} There is a constant $C$ such that for all sufficiently large primes $p$, all $b\not\equiv0\pmod p$ and all $N<p$ with $\log N\ge\frac12(\log p)^{9/10}$,
\[
\Bigl|\sum_{n\le N}e\Bigl(\frac{bn^*}p\Bigr)\Bigr|\le C(\log\log p)^3(\log p)(\log N)^{-3/2}N .
\]
This is exactly Lemma~\ref{lem:bg}, in the same range. The Bourgain--Garaev bound~\cite[Theorem~16]{b8} is not formalized. The weaker hypothesis \texttt{HKS} asks for some $\delta>0$, $\theta<1$ and $C$ such that for large $p$, all $b\not\equiv0\pmod p$ and all $N<p$ with $\log N\ge(\log p)^\theta$,
\[
\Bigl|\sum_{n\le N}e(bn^*/p)\Bigr|\le CN(\log p)^{-\delta}.
\]
Lean proves \lean{ThetaSum.HKS_of_HBG}: (HBG) implies (HKS), with $\theta=\frac9{10}$ and $\delta=\frac14$.

\subsection*{The main formal theorem}
The theorem \lean{ThetaSum.conjecture_3_3_of_BG} (file \lean{ThetaSum/Headline.lean}) reads
\begin{leancode}
theorem conjecture_3_3_of_BG (hBG : HBG) : $\forall$ C$\sb{1}$ C$\sb{2}$ $\varepsilon$ : $\R$, 0 < C$\sb{1}$ $\to$ 0 < C$\sb{2}$ $\to$ 0 < $\varepsilon$ $\to$
    $\forall\sp{f}$ k : $\N$ in Filter.atTop, k.Prime $\to$
      C$\sb{1}$ * k < (S k : $\R$) $\wedge$ (S k : $\R$) < C$\sb{2}$ * (k : $\R$) ^ (1 + $\varepsilon$)
\end{leancode}
It proves Conjecture~3.3 (both halves) for the definition of $S$ in~\cite{b11}, assuming only (HBG). The same hypothesis gives the statement \lean{conjecture_4_4} of~\cite{b11} (Conjecture~4.4 of the first arXiv version: $S(k)>nk$ for each $n\in\N$ and all large primes $k$, which is equivalent to the lower half of Conjecture~3.3; \lean{conjecture_4_4_of_BG}), and Conjectures~1.1, 3.1 and~3.2 for all large $p$ (\lean{conjecture_1_1_eventually_of_BG}, \lean{conjecture_3_1_eventually_of_BG}, \lean{conjecture_3_2_eventually_of_BG}). The file \lean{ThetaSum/WeakHyp.lean} contains \lean{conjecture_3_3_of_HKS}, the same statement from (HKS).

\subsection*{Numbering}
The formal-conjectures file uses the numbering of the first arXiv version of~\cite{b1}: Conjecture~3.1 is \lean{conjecture_4_2}, Conjecture~3.2 is \lean{conjecture_4_3}, and \lean{conjecture_4_1} is the inequality $S(k)>k$ for primes $k>5$.

\subsection*{What else is formalized}
Without any hypothesis, Lean proves: Proposition~\ref{prop:upper} (\lean{theorem_A}) and the upper half of Conjecture~3.3 (\lean{conjecture_3_3_upper}); Theorem~\ref{thm:decomp}(i), where the identity is \lean{TheoremB.S_eq_Psi_Phi} and the asymptotic formula for $\Psi(p)$, with error at most $\varepsilon p$, is \lean{TheoremB.Psi_asymp}; Theorem~\ref{thm:decomp}(ii) in two forms, both for $K=\lfloor p^{3/4}(\log p)^4\rfloor$: \lean{ThetaSum.hBweak34}, with an unspecified constant $c$ in place of $c_0$ and error $O(p)$, and \lean{ThetaSum.theorem_B_c0}, with the constant $c_0$ and error at most $\varepsilon p$; Lemma~\ref{lem:kloosterman} (\lean{ThetaSum.hK34}); Theorem~\ref{thm:divisors} (\lean{ThetaSum.theorem_D}); Lemma~\ref{lem:counting} (\lean{ThetaSum.S_eq_lemma25}, with \lean{Tq_odd}, \lean{Ffl_add} and \lean{Ffl_recip}); and the finite checks for $p\le2\cdot10^4$ (\lean{S_pos_upto}, \lean{S_gt_self_upto}, \lean{S_gt_two_mul_upto}, \lean{S_gt_three_mul_upto}, \lean{S_233}, \lean{S_3119}), computed by the Lean kernel.

The formalization also contains \lean{ThetaSum.theorem_C_star}, which gives the error term $p(\log p)^{2/3}(\log\log p)^{17/3}$ under the stronger hypothesis \lean{HBGfull} (the Bourgain--Garaev bound for $\log N\ge A(\log p)^{2/3}(\log\log p)^{8/3}$), and \lean{ThetaSum.theorem_D_star}, the variant of Theorem~\ref{thm:divisors} with $\log T\ge(\log\log q)^3$, stated for $\log\log q\ge1000$. These results are not used in this paper.

The formal proofs differ from the paper in these respects. Some constants differ. Lemma~\ref{lem:generic} is formalized in an elementary version that does not use Vaaler's lemma. The reciprocity law (\lean{Recip.reciprocity}) is proved without contour integration. The Lean assembly proves $S(p)/p\to\infty$ (that is, $S(p)=\frac p2\log p+o(p\log p)$) rather than the explicit exponent $\frac{19}{20}$ of Theorem~\ref{thm:main}.

\subsection*{How to check}
In the directory \texttt{lean} of the repository, run
\begin{leancode}
lake exe cache get
lake build
lake env lean ThetaSum/MainTheorems.lean
\end{leancode}
The first command downloads the compiled Mathlib, the second builds the project (about 30 minutes on four cores), and the third prints the statement and the axioms of each headline result. The hypothesis \texttt{HBG} appears as an argument of the theorem, not as an axiom.

\end{document}
